% Options for packages loaded elsewhere
\PassOptionsToPackage{unicode}{hyperref}
\PassOptionsToPackage{hyphens}{url}
\documentclass[
]{article}
\usepackage{xcolor}
\usepackage[margin=1in]{geometry}
\usepackage{amsmath,amssymb}
\setcounter{secnumdepth}{-\maxdimen} % remove section numbering
\usepackage{iftex}
\ifPDFTeX
  \usepackage[T1]{fontenc}
  \usepackage[utf8]{inputenc}
  \usepackage{textcomp} % provide euro and other symbols
\else % if luatex or xetex
  \usepackage{unicode-math} % this also loads fontspec
  \defaultfontfeatures{Scale=MatchLowercase}
  \defaultfontfeatures[\rmfamily]{Ligatures=TeX,Scale=1}
\fi
\usepackage{lmodern}
\ifPDFTeX\else
  % xetex/luatex font selection
\fi
% Use upquote if available, for straight quotes in verbatim environments
\IfFileExists{upquote.sty}{\usepackage{upquote}}{}
\IfFileExists{microtype.sty}{% use microtype if available
  \usepackage[]{microtype}
  \UseMicrotypeSet[protrusion]{basicmath} % disable protrusion for tt fonts
}{}
\makeatletter
\@ifundefined{KOMAClassName}{% if non-KOMA class
  \IfFileExists{parskip.sty}{%
    \usepackage{parskip}
  }{% else
    \setlength{\parindent}{0pt}
    \setlength{\parskip}{6pt plus 2pt minus 1pt}}
}{% if KOMA class
  \KOMAoptions{parskip=half}}
\makeatother
\usepackage{color}
\usepackage{fancyvrb}
\newcommand{\VerbBar}{|}
\newcommand{\VERB}{\Verb[commandchars=\\\{\}]}
\DefineVerbatimEnvironment{Highlighting}{Verbatim}{commandchars=\\\{\}}
% Add ',fontsize=\small' for more characters per line
\usepackage{framed}
\definecolor{shadecolor}{RGB}{248,248,248}
\newenvironment{Shaded}{\begin{snugshade}}{\end{snugshade}}
\newcommand{\AlertTok}[1]{\textcolor[rgb]{0.94,0.16,0.16}{#1}}
\newcommand{\AnnotationTok}[1]{\textcolor[rgb]{0.56,0.35,0.01}{\textbf{\textit{#1}}}}
\newcommand{\AttributeTok}[1]{\textcolor[rgb]{0.13,0.29,0.53}{#1}}
\newcommand{\BaseNTok}[1]{\textcolor[rgb]{0.00,0.00,0.81}{#1}}
\newcommand{\BuiltInTok}[1]{#1}
\newcommand{\CharTok}[1]{\textcolor[rgb]{0.31,0.60,0.02}{#1}}
\newcommand{\CommentTok}[1]{\textcolor[rgb]{0.56,0.35,0.01}{\textit{#1}}}
\newcommand{\CommentVarTok}[1]{\textcolor[rgb]{0.56,0.35,0.01}{\textbf{\textit{#1}}}}
\newcommand{\ConstantTok}[1]{\textcolor[rgb]{0.56,0.35,0.01}{#1}}
\newcommand{\ControlFlowTok}[1]{\textcolor[rgb]{0.13,0.29,0.53}{\textbf{#1}}}
\newcommand{\DataTypeTok}[1]{\textcolor[rgb]{0.13,0.29,0.53}{#1}}
\newcommand{\DecValTok}[1]{\textcolor[rgb]{0.00,0.00,0.81}{#1}}
\newcommand{\DocumentationTok}[1]{\textcolor[rgb]{0.56,0.35,0.01}{\textbf{\textit{#1}}}}
\newcommand{\ErrorTok}[1]{\textcolor[rgb]{0.64,0.00,0.00}{\textbf{#1}}}
\newcommand{\ExtensionTok}[1]{#1}
\newcommand{\FloatTok}[1]{\textcolor[rgb]{0.00,0.00,0.81}{#1}}
\newcommand{\FunctionTok}[1]{\textcolor[rgb]{0.13,0.29,0.53}{\textbf{#1}}}
\newcommand{\ImportTok}[1]{#1}
\newcommand{\InformationTok}[1]{\textcolor[rgb]{0.56,0.35,0.01}{\textbf{\textit{#1}}}}
\newcommand{\KeywordTok}[1]{\textcolor[rgb]{0.13,0.29,0.53}{\textbf{#1}}}
\newcommand{\NormalTok}[1]{#1}
\newcommand{\OperatorTok}[1]{\textcolor[rgb]{0.81,0.36,0.00}{\textbf{#1}}}
\newcommand{\OtherTok}[1]{\textcolor[rgb]{0.56,0.35,0.01}{#1}}
\newcommand{\PreprocessorTok}[1]{\textcolor[rgb]{0.56,0.35,0.01}{\textit{#1}}}
\newcommand{\RegionMarkerTok}[1]{#1}
\newcommand{\SpecialCharTok}[1]{\textcolor[rgb]{0.81,0.36,0.00}{\textbf{#1}}}
\newcommand{\SpecialStringTok}[1]{\textcolor[rgb]{0.31,0.60,0.02}{#1}}
\newcommand{\StringTok}[1]{\textcolor[rgb]{0.31,0.60,0.02}{#1}}
\newcommand{\VariableTok}[1]{\textcolor[rgb]{0.00,0.00,0.00}{#1}}
\newcommand{\VerbatimStringTok}[1]{\textcolor[rgb]{0.31,0.60,0.02}{#1}}
\newcommand{\WarningTok}[1]{\textcolor[rgb]{0.56,0.35,0.01}{\textbf{\textit{#1}}}}
\usepackage{longtable,booktabs,array}
\newcounter{none} % for unnumbered tables
\usepackage{calc} % for calculating minipage widths
% Correct order of tables after \paragraph or \subparagraph
\usepackage{etoolbox}
\makeatletter
\patchcmd\longtable{\par}{\if@noskipsec\mbox{}\fi\par}{}{}
\makeatother
% Allow footnotes in longtable head/foot
\IfFileExists{footnotehyper.sty}{\usepackage{footnotehyper}}{\usepackage{footnote}}
\makesavenoteenv{longtable}
\usepackage{graphicx}
\makeatletter
\newsavebox\pandoc@box
\newcommand*\pandocbounded[1]{% scales image to fit in text height/width
  \sbox\pandoc@box{#1}%
  \Gscale@div\@tempa{\textheight}{\dimexpr\ht\pandoc@box+\dp\pandoc@box\relax}%
  \Gscale@div\@tempb{\linewidth}{\wd\pandoc@box}%
  \ifdim\@tempb\p@<\@tempa\p@\let\@tempa\@tempb\fi% select the smaller of both
  \ifdim\@tempa\p@<\p@\scalebox{\@tempa}{\usebox\pandoc@box}%
  \else\usebox{\pandoc@box}%
  \fi%
}
% Set default figure placement to htbp
\def\fps@figure{htbp}
\makeatother
\setlength{\emergencystretch}{3em} % prevent overfull lines
\providecommand{\tightlist}{%
  \setlength{\itemsep}{0pt}\setlength{\parskip}{0pt}}
\usepackage{bookmark}
\IfFileExists{xurl.sty}{\usepackage{xurl}}{} % add URL line breaks if available
\urlstyle{same}
\hypersetup{
  pdftitle={What Kind of Payment Network Are We?},
  pdfauthor={Team ERROR 404 \textbar{} Head of Analytics, Elo \textbar{} Prepared for the Q4 Strategy Offsite},
  hidelinks,
  pdfcreator={LaTeX via pandoc}}

\title{What Kind of Payment Network Are We?}
\usepackage{etoolbox}
\makeatletter
\providecommand{\subtitle}[1]{% add subtitle to \maketitle
  \apptocmd{\@title}{\par {\large #1 \par}}{}{}
}
\makeatother
\subtitle{A portfolio-level review of the Elo card ledger, merchant base
and promotions programme}
\author{Team ERROR 404 \textbar{} Head of Analytics, Elo \textbar{}
Prepared for the Q4 Strategy Offsite}
\date{28 September 2026}

\begin{document}
\maketitle

{
\setcounter{tocdepth}{3}
\tableofcontents
}
\section{The brief, and how we chose to answer
it}\label{the-brief-and-how-we-chose-to-answer-it}

The Board asked one question: \textbf{based on everything in this data,
what kind of payment network are we, and what should that tell us about
where to point the promotions programme in Quarter 4?}

That is deliberately not a modelling question. Nobody is asking for an
uplift score. They are asking for an identity --- and then for the
operating consequence of that identity. So we structured the work as a
chain of five questions, each one earning the right to ask the next:

{\def\LTcaptype{none} % do not increment counter
\begin{longtable}[]{@{}
  >{\raggedright\arraybackslash}p{(\linewidth - 4\tabcolsep) * \real{0.0769}}
  >{\raggedright\arraybackslash}p{(\linewidth - 4\tabcolsep) * \real{0.2564}}
  >{\raggedright\arraybackslash}p{(\linewidth - 4\tabcolsep) * \real{0.6667}}@{}}
\toprule\noalign{}
\begin{minipage}[b]{\linewidth}\raggedright
\#
\end{minipage} & \begin{minipage}[b]{\linewidth}\raggedright
Question
\end{minipage} & \begin{minipage}[b]{\linewidth}\raggedright
Why it has to come first
\end{minipage} \\
\midrule\noalign{}
\endhead
\bottomrule\noalign{}
\endlastfoot
1 & Can we even read the money? & \texttt{purchase\_amount} is
normalised. Without real currency no finding is presentable to a
Board. \\
2 & Who transacts on our cards? & Concentration decides whether
promotions are a mass-market or a targeted instrument. \\
3 & How is the network actually used? & Ticket size and instalment
behaviour \emph{are} the identity answer. \\
4 & What does the merchant side look like? & Promotions are paid to
merchants. Their size and health decide the price. \\
5 & Where does discovery happen? & The new-merchant window is the only
place in this export where a \emph{new} commercial relationship is
visible. That is what a promotion is trying to create. \\
\end{longtable}
}

The recommendation in Section 9 falls directly out of question 5 crossed
with question 4.

\begin{center}\rule{0.5\linewidth}{0.5pt}\end{center}

\section{Data assembly and what we found wrong with
it}\label{data-assembly-and-what-we-found-wrong-with-it}

Three files were handed over. Before any analysis, we audited them.

\begin{Shaded}
\begin{Highlighting}[]
\NormalTok{hist\_raw  }\OtherTok{\textless{}{-}} \FunctionTok{read\_fast}\NormalTok{(}\FunctionTok{file.path}\NormalTok{(DATA\_DIR, }\StringTok{"historical\_transactions\_100k.csv"}\NormalTok{))}
\NormalTok{new\_raw   }\OtherTok{\textless{}{-}} \FunctionTok{read\_fast}\NormalTok{(}\FunctionTok{file.path}\NormalTok{(DATA\_DIR, }\StringTok{"new\_merchant\_transactions.csv"}\NormalTok{))}
\NormalTok{merch\_raw }\OtherTok{\textless{}{-}} \FunctionTok{read\_fast}\NormalTok{(}\FunctionTok{file.path}\NormalTok{(DATA\_DIR, }\StringTok{"merchants1.csv"}\NormalTok{))}

\FunctionTok{data.frame}\NormalTok{(}
  \AttributeTok{File =} \FunctionTok{c}\NormalTok{(}\StringTok{"historical\_transactions\_100k.csv"}\NormalTok{, }\StringTok{"new\_merchant\_transactions.csv"}\NormalTok{, }\StringTok{"merchants1.csv"}\NormalTok{),}
  \StringTok{\textasciigrave{}}\AttributeTok{Rows as delivered}\StringTok{\textasciigrave{}} \OtherTok{=} \FunctionTok{c}\NormalTok{(}\FunctionTok{nrow}\NormalTok{(hist\_raw), }\FunctionTok{nrow}\NormalTok{(new\_raw), }\FunctionTok{nrow}\NormalTok{(merch\_raw)),}
  \AttributeTok{Columns =} \FunctionTok{c}\NormalTok{(}\FunctionTok{ncol}\NormalTok{(hist\_raw), }\FunctionTok{ncol}\NormalTok{(new\_raw), }\FunctionTok{ncol}\NormalTok{(merch\_raw)),}
  \AttributeTok{check.names =} \ConstantTok{FALSE}\NormalTok{) }\SpecialCharTok{\%\textgreater{}\%} \FunctionTok{kable}\NormalTok{(}\AttributeTok{format.args =} \FunctionTok{list}\NormalTok{(}\AttributeTok{big.mark =} \StringTok{","}\NormalTok{))}
\end{Highlighting}
\end{Shaded}

{\def\LTcaptype{none} % do not increment counter
\begin{longtable}[]{@{}lrr@{}}
\toprule\noalign{}
File & Rows as delivered & Columns \\
\midrule\noalign{}
\endhead
\bottomrule\noalign{}
\endlastfoot
historical\_transactions\_100k.csv & 100,000 & 14 \\
new\_merchant\_transactions.csv & 1,048,575 & 14 \\
merchants1.csv & 334,696 & 22 \\
\end{longtable}
}

\subsection{Three defects worth
naming}\label{three-defects-worth-naming}

\textbf{One --- the new-merchant file is 90\% empty.} It arrives with
1,048,575 rows, which is exactly the Excel worksheet limit minus the
header. Everything past row 100,000 is a row of commas. Any team that
reported ``1.05 million new-merchant transactions'' to the Board would
have been wrong by a factor of ten.

\begin{Shaded}
\begin{Highlighting}[]
\NormalTok{filler }\OtherTok{\textless{}{-}} \FunctionTok{sum}\NormalTok{(}\FunctionTok{is.na}\NormalTok{(new\_raw}\SpecialCharTok{$}\NormalTok{card\_id) }\SpecialCharTok{|}\NormalTok{ new\_raw}\SpecialCharTok{$}\NormalTok{card\_id }\SpecialCharTok{==} \StringTok{""}\NormalTok{)}
\NormalTok{new\_raw }\OtherTok{\textless{}{-}}\NormalTok{ new\_raw[}\SpecialCharTok{!}\FunctionTok{is.na}\NormalTok{(new\_raw}\SpecialCharTok{$}\NormalTok{card\_id) }\SpecialCharTok{\&}\NormalTok{ new\_raw}\SpecialCharTok{$}\NormalTok{card\_id }\SpecialCharTok{!=} \StringTok{""}\NormalTok{, ]}
\FunctionTok{cat}\NormalTok{(}\StringTok{"Empty filler rows removed:"}\NormalTok{, }\FunctionTok{format}\NormalTok{(filler, }\AttributeTok{big.mark =} \StringTok{","}\NormalTok{),}
    \StringTok{"}\SpecialCharTok{\textbackslash{}n}\StringTok{Real new{-}merchant transactions:"}\NormalTok{, }\FunctionTok{format}\NormalTok{(}\FunctionTok{nrow}\NormalTok{(new\_raw), }\AttributeTok{big.mark =} \StringTok{","}\NormalTok{))}
\end{Highlighting}
\end{Shaded}

\begin{verbatim}
## Empty filler rows removed: 948,575 
## Real new-merchant transactions: 100,000
\end{verbatim}

\textbf{Two --- the merchant reference file repeats itself.} 63
\texttt{merchant\_id} values appear twice. Left alone they double-count
in every supply-side share.

\textbf{Three --- the growth ratios contain literal infinities.} Three
merchants had zero transactions in their base month, so
\texttt{avg\_purchases\_lag*} divides by zero.

\begin{Shaded}
\begin{Highlighting}[]
\NormalTok{dup\_merch }\OtherTok{\textless{}{-}} \FunctionTok{sum}\NormalTok{(}\FunctionTok{duplicated}\NormalTok{(merch\_raw}\SpecialCharTok{$}\NormalTok{merchant\_id))}
\NormalTok{merch }\OtherTok{\textless{}{-}}\NormalTok{ merch\_raw[}\SpecialCharTok{!}\FunctionTok{duplicated}\NormalTok{(merch\_raw}\SpecialCharTok{$}\NormalTok{merchant\_id), ]}

\NormalTok{lag\_cols }\OtherTok{\textless{}{-}} \FunctionTok{c}\NormalTok{(}\StringTok{"avg\_sales\_lag3"}\NormalTok{, }\StringTok{"avg\_purchases\_lag3"}\NormalTok{, }\StringTok{"avg\_sales\_lag6"}\NormalTok{,}
              \StringTok{"avg\_purchases\_lag6"}\NormalTok{, }\StringTok{"avg\_sales\_lag12"}\NormalTok{, }\StringTok{"avg\_purchases\_lag12"}\NormalTok{)}
\ControlFlowTok{for}\NormalTok{ (cc }\ControlFlowTok{in}\NormalTok{ lag\_cols) merch[[cc]] }\OtherTok{\textless{}{-}} \FunctionTok{suppressWarnings}\NormalTok{(}\FunctionTok{as.numeric}\NormalTok{(merch[[cc]]))}
\ControlFlowTok{for}\NormalTok{ (cc }\ControlFlowTok{in} \FunctionTok{c}\NormalTok{(}\StringTok{"subsector\_id"}\NormalTok{, }\StringTok{"state\_id"}\NormalTok{, }\StringTok{"city\_id"}\NormalTok{, }\StringTok{"merchant\_category\_id"}\NormalTok{,}
             \StringTok{"numerical\_1"}\NormalTok{, }\StringTok{"numerical\_2"}\NormalTok{, }\StringTok{"category\_2"}\NormalTok{,}
             \StringTok{"active\_months\_lag3"}\NormalTok{, }\StringTok{"active\_months\_lag6"}\NormalTok{, }\StringTok{"active\_months\_lag12"}\NormalTok{))}
\NormalTok{  merch[[cc]] }\OtherTok{\textless{}{-}} \FunctionTok{suppressWarnings}\NormalTok{(}\FunctionTok{as.numeric}\NormalTok{(merch[[cc]]))}

\FunctionTok{cat}\NormalTok{(}\StringTok{"Duplicate merchant rows dropped:"}\NormalTok{, dup\_merch,}
    \StringTok{"}\SpecialCharTok{\textbackslash{}n}\StringTok{Unique merchants:"}\NormalTok{, }\FunctionTok{format}\NormalTok{(}\FunctionTok{nrow}\NormalTok{(merch), }\AttributeTok{big.mark =} \StringTok{","}\NormalTok{),}
    \StringTok{"}\SpecialCharTok{\textbackslash{}n}\StringTok{Non{-}finite growth ratios:"}\NormalTok{, }\FunctionTok{sum}\NormalTok{(}\SpecialCharTok{!}\FunctionTok{is.finite}\NormalTok{(merch}\SpecialCharTok{$}\NormalTok{avg\_purchases\_lag3)))}
\end{Highlighting}
\end{Shaded}

\begin{verbatim}
## Duplicate merchant rows dropped: 63 
## Unique merchants: 334,633 
## Non-finite growth ratios: 3
\end{verbatim}

\subsection{A structural limit we have to be honest
about}\label{a-structural-limit-we-have-to-be-honest-about}

The two transaction files were sampled independently. The historical
ledger covers \textbf{463 cards} in depth; the new-merchant window
covers \textbf{7,530 cards} shallowly; only \textbf{16 cards appear in
both}.

\begin{Shaded}
\begin{Highlighting}[]
\FunctionTok{cat}\NormalTok{(}\StringTok{"Cards in historical ledger :"}\NormalTok{, }\FunctionTok{n\_distinct}\NormalTok{(hist\_raw}\SpecialCharTok{$}\NormalTok{card\_id),}
    \StringTok{"}\SpecialCharTok{\textbackslash{}n}\StringTok{Cards in novelty window    :"}\NormalTok{, }\FunctionTok{n\_distinct}\NormalTok{(new\_raw}\SpecialCharTok{$}\NormalTok{card\_id),}
    \StringTok{"}\SpecialCharTok{\textbackslash{}n}\StringTok{Cards present in both      :"}\NormalTok{,}
    \FunctionTok{length}\NormalTok{(}\FunctionTok{intersect}\NormalTok{(}\FunctionTok{unique}\NormalTok{(hist\_raw}\SpecialCharTok{$}\NormalTok{card\_id), }\FunctionTok{unique}\NormalTok{(new\_raw}\SpecialCharTok{$}\NormalTok{card\_id))))}
\end{Highlighting}
\end{Shaded}

\begin{verbatim}
## Cards in historical ledger : 463 
## Cards in novelty window    : 7530 
## Cards present in both      : 16
\end{verbatim}

So we \textbf{cannot} ask ``did card X change behaviour when it met a
new merchant''. We can --- and do --- ask the portfolio-level version:
\emph{how does spending at merchants a card already knows differ from
spending at merchants it has never visited?} Every novelty finding in
this document is a comparison of \textbf{two windows of the same
network}, not a within-card before/after. Stated up front so no one
over-reads it.

The card-level file described in the data dictionary
(\texttt{train.csv}, carrying \texttt{first\_active\_month} and the
loyalty \texttt{target}) was not part of the export. We therefore
rebuild tenure and value from observed behaviour rather than using the
supplied loyalty score. Section 10 sets out what that costs us.

\begin{center}\rule{0.5\linewidth}{0.5pt}\end{center}

\section{\texorpdfstring{Reading the money: de-normalising
\texttt{purchase\_amount}}{Reading the money: de-normalising purchase\_amount}}\label{reading-the-money-de-normalising-purchase_amount}

\texttt{purchase\_amount} is a normalised quantity ranging from about
−0.75 to +11,270. Nobody can take that to a Board. It is, however, an
\textbf{affine transform} of the real Brazilian Real amount, and an
affine transform can be recovered.

Two anchors do it. The floor of the distribution must correspond to the
smallest possible real purchase, and the bulk of a retail ledger must
land in sensible everyday money. Solving on that basis gives:

\[\text{BRL} = \frac{\texttt{purchase\_amount}}{0.00150265} + 497.06\]

\begin{Shaded}
\begin{Highlighting}[]
\NormalTok{SCALE }\OtherTok{\textless{}{-}} \FloatTok{0.00150265}\NormalTok{; OFFSET }\OtherTok{\textless{}{-}} \FloatTok{497.06}
\NormalTok{to\_brl }\OtherTok{\textless{}{-}} \ControlFlowTok{function}\NormalTok{(x) x }\SpecialCharTok{/}\NormalTok{ SCALE }\SpecialCharTok{+}\NormalTok{ OFFSET}

\NormalTok{prep\_txn }\OtherTok{\textless{}{-}} \ControlFlowTok{function}\NormalTok{(df) \{}
\NormalTok{  num\_cols }\OtherTok{\textless{}{-}} \FunctionTok{c}\NormalTok{(}\StringTok{"purchase\_amount"}\NormalTok{, }\StringTok{"installments"}\NormalTok{, }\StringTok{"month\_lag"}\NormalTok{, }\StringTok{"city\_id"}\NormalTok{,}
                \StringTok{"state\_id"}\NormalTok{, }\StringTok{"subsector\_id"}\NormalTok{, }\StringTok{"merchant\_category\_id"}\NormalTok{, }\StringTok{"category\_2"}\NormalTok{)}
  \ControlFlowTok{for}\NormalTok{ (cc }\ControlFlowTok{in} \FunctionTok{intersect}\NormalTok{(num\_cols, }\FunctionTok{names}\NormalTok{(df)))}
\NormalTok{    df[[cc]] }\OtherTok{\textless{}{-}} \FunctionTok{suppressWarnings}\NormalTok{(}\FunctionTok{as.numeric}\NormalTok{(df[[cc]]))}
\NormalTok{  df }\SpecialCharTok{\%\textgreater{}\%} \FunctionTok{mutate}\NormalTok{(}
    \AttributeTok{amount\_brl  =} \FunctionTok{to\_brl}\NormalTok{(purchase\_amount),}
    \AttributeTok{purchase\_ts =} \FunctionTok{as.POSIXct}\NormalTok{(purchase\_date, }\AttributeTok{format =} \StringTok{"\%d{-}\%m{-}\%Y \%H:\%M"}\NormalTok{, }\AttributeTok{tz =} \StringTok{"UTC"}\NormalTok{),}
    \AttributeTok{month       =} \FunctionTok{format}\NormalTok{(purchase\_ts, }\StringTok{"\%Y{-}\%m"}\NormalTok{),}
    \AttributeTok{hour        =} \FunctionTok{as.integer}\NormalTok{(}\FunctionTok{format}\NormalTok{(purchase\_ts, }\StringTok{"\%H"}\NormalTok{)),}
    \AttributeTok{instal      =} \FunctionTok{ifelse}\NormalTok{(}\FunctionTok{is.na}\NormalTok{(installments) }\SpecialCharTok{|}\NormalTok{ installments }\SpecialCharTok{\textless{}} \DecValTok{0}\NormalTok{, }\ConstantTok{NA\_real\_}\NormalTok{,}
                         \FunctionTok{as.numeric}\NormalTok{(installments)),}
    \AttributeTok{instal\_band =} \FunctionTok{cut}\NormalTok{(instal, }\FunctionTok{c}\NormalTok{(}\SpecialCharTok{{-}}\DecValTok{1}\NormalTok{, }\DecValTok{0}\NormalTok{, }\DecValTok{1}\NormalTok{, }\DecValTok{3}\NormalTok{, }\DecValTok{6}\NormalTok{, }\DecValTok{12}\NormalTok{, }\DecValTok{999}\NormalTok{),}
                      \AttributeTok{labels =} \FunctionTok{c}\NormalTok{(}\StringTok{"0 (lump sum)"}\NormalTok{, }\StringTok{"1 (single)"}\NormalTok{, }\StringTok{"2{-}3"}\NormalTok{, }\StringTok{"4{-}6"}\NormalTok{, }\StringTok{"7{-}12"}\NormalTok{, }\StringTok{"13+"}\NormalTok{)),}
    \AttributeTok{approved    =}\NormalTok{ authorized\_flag }\SpecialCharTok{==} \StringTok{"Y"}\NormalTok{)}
\NormalTok{\}}
\NormalTok{hist\_txn }\OtherTok{\textless{}{-}} \FunctionTok{prep\_txn}\NormalTok{(hist\_raw)}
\NormalTok{new\_txn  }\OtherTok{\textless{}{-}} \FunctionTok{prep\_txn}\NormalTok{(new\_raw)}

\FunctionTok{round}\NormalTok{(}\FunctionTok{quantile}\NormalTok{(hist\_txn}\SpecialCharTok{$}\NormalTok{amount\_brl, }\FunctionTok{c}\NormalTok{(}\DecValTok{0}\NormalTok{, .}\DecValTok{01}\NormalTok{, .}\DecValTok{25}\NormalTok{, .}\DecValTok{5}\NormalTok{, .}\DecValTok{75}\NormalTok{, .}\DecValTok{95}\NormalTok{, .}\DecValTok{99}\NormalTok{, }\DecValTok{1}\NormalTok{)), }\DecValTok{2}\NormalTok{) }\SpecialCharTok{\%\textgreater{}\%}
  \FunctionTok{as.data.frame}\NormalTok{() }\SpecialCharTok{\%\textgreater{}\%} \FunctionTok{setNames}\NormalTok{(}\StringTok{"Amount (BRL)"}\NormalTok{) }\SpecialCharTok{\%\textgreater{}\%} \FunctionTok{kable}\NormalTok{(}\AttributeTok{format.args =} \FunctionTok{list}\NormalTok{(}\AttributeTok{big.mark =} \StringTok{","}\NormalTok{))}
\end{Highlighting}
\end{Shaded}

{\def\LTcaptype{none} % do not increment counter
\begin{longtable}[]{@{}lr@{}}
\toprule\noalign{}
& Amount (BRL) \\
\midrule\noalign{}
\endhead
\bottomrule\noalign{}
\endlastfoot
0\% & 0.01 \\
1\% & 2.78 \\
25\% & 16.32 \\
50\% & 34.90 \\
75\% & 82.62 \\
95\% & 300.00 \\
99\% & 1,030.80 \\
100\% & 7,500,358.32 \\
\end{longtable}
}

\textbf{The validation is the distribution itself.} The minimum lands at
R\$0.01, the median at roughly \textbf{R\$35}, the 75th percentile at
\textbf{R\$83} and the 99th at about \textbf{R\$1,031}. Not one
transaction is negative. A wrong constant would have produced negative
money at the bottom of the ledger, and it does not. From here on, every
currency figure in this report is real Reais.

\subsection{Trimming the tail}\label{trimming-the-tail}

99 records --- one in a thousand --- reach into the millions of Reais
and are almost certainly settlement or reversal artefacts rather than
shopping. One of them alone would move the portfolio average ticket by
more than every finding in this report combined. We exclude the top
0.1\% from all value statistics and say so wherever it matters.

\begin{Shaded}
\begin{Highlighting}[]
\NormalTok{CAP    }\OtherTok{\textless{}{-}} \FunctionTok{quantile}\NormalTok{(hist\_txn}\SpecialCharTok{$}\NormalTok{amount\_brl, }\FloatTok{0.999}\NormalTok{)}
\NormalTok{hist\_w }\OtherTok{\textless{}{-}}\NormalTok{ hist\_txn }\SpecialCharTok{\%\textgreater{}\%} \FunctionTok{filter}\NormalTok{(amount\_brl }\SpecialCharTok{\textless{}=}\NormalTok{ CAP)}
\NormalTok{new\_w  }\OtherTok{\textless{}{-}}\NormalTok{ new\_txn  }\SpecialCharTok{\%\textgreater{}\%} \FunctionTok{filter}\NormalTok{(amount\_brl }\SpecialCharTok{\textless{}=} \FunctionTok{quantile}\NormalTok{(amount\_brl, }\FloatTok{0.999}\NormalTok{))}
\NormalTok{hist\_ok }\OtherTok{\textless{}{-}}\NormalTok{ hist\_w }\SpecialCharTok{\%\textgreater{}\%} \FunctionTok{filter}\NormalTok{(approved)}

\FunctionTok{cat}\NormalTok{(}\StringTok{"Trim threshold:"}\NormalTok{, }\FunctionTok{comma}\NormalTok{(}\FunctionTok{round}\NormalTok{(CAP)), }\StringTok{"BRL | records excluded:"}\NormalTok{,}
    \FunctionTok{sum}\NormalTok{(hist\_txn}\SpecialCharTok{$}\NormalTok{amount\_brl }\SpecialCharTok{\textgreater{}}\NormalTok{ CAP),}
    \StringTok{"}\SpecialCharTok{\textbackslash{}n}\StringTok{Mean ticket before trim: R$"}\NormalTok{, }\FunctionTok{round}\NormalTok{(}\FunctionTok{mean}\NormalTok{(hist\_txn}\SpecialCharTok{$}\NormalTok{amount\_brl), }\DecValTok{2}\NormalTok{),}
    \StringTok{" | after trim: R$"}\NormalTok{, }\FunctionTok{round}\NormalTok{(}\FunctionTok{mean}\NormalTok{(hist\_w}\SpecialCharTok{$}\NormalTok{amount\_brl), }\DecValTok{2}\NormalTok{))}
\end{Highlighting}
\end{Shaded}

\begin{verbatim}
## Trim threshold: 4,000 BRL | records excluded: 99 
## Mean ticket before trim: R$ 272.11  | after trim: R$ 88.24
\end{verbatim}

The average ticket falls from R\$272 to \textbf{R\$88} once 99 records
are removed. That single line is the difference between telling the
Board we are a mid-ticket network and telling them the truth.

\begin{center}\rule{0.5\linewidth}{0.5pt}\end{center}

\section{Who transacts on our cards?}\label{who-transacts-on-our-cards}

\begin{Shaded}
\begin{Highlighting}[]
\NormalTok{ref\_date }\OtherTok{\textless{}{-}} \FunctionTok{max}\NormalTok{(hist\_ok}\SpecialCharTok{$}\NormalTok{purchase\_ts, }\AttributeTok{na.rm =} \ConstantTok{TRUE}\NormalTok{)}

\NormalTok{card }\OtherTok{\textless{}{-}}\NormalTok{ hist\_ok }\SpecialCharTok{\%\textgreater{}\%}
  \FunctionTok{group\_by}\NormalTok{(card\_id) }\SpecialCharTok{\%\textgreater{}\%}
  \FunctionTok{summarise}\NormalTok{(}\AttributeTok{txns =} \FunctionTok{n}\NormalTok{(), }\AttributeTok{spend =} \FunctionTok{sum}\NormalTok{(amount\_brl), }\AttributeTok{avg\_ticket =} \FunctionTok{mean}\NormalTok{(amount\_brl),}
            \AttributeTok{merchants =} \FunctionTok{n\_distinct}\NormalTok{(merchant\_id), }\AttributeTok{subsectors =} \FunctionTok{n\_distinct}\NormalTok{(subsector\_id),}
            \AttributeTok{first\_seen =} \FunctionTok{min}\NormalTok{(purchase\_ts), }\AttributeTok{last\_seen =} \FunctionTok{max}\NormalTok{(purchase\_ts), }\AttributeTok{.groups =} \StringTok{"drop"}\NormalTok{) }\SpecialCharTok{\%\textgreater{}\%}
  \FunctionTok{mutate}\NormalTok{(}\AttributeTok{recency\_days  =} \FunctionTok{as.numeric}\NormalTok{(}\FunctionTok{difftime}\NormalTok{(ref\_date, last\_seen, }\AttributeTok{units =} \StringTok{"days"}\NormalTok{)),}
         \AttributeTok{active\_months =} \FunctionTok{pmax}\NormalTok{(}\FunctionTok{as.numeric}\NormalTok{(}\FunctionTok{difftime}\NormalTok{(last\_seen, first\_seen, }\AttributeTok{units =} \StringTok{"days"}\NormalTok{)) }\SpecialCharTok{/} \FloatTok{30.44}\NormalTok{, }\DecValTok{1}\NormalTok{),}
         \AttributeTok{txns\_pm  =}\NormalTok{ txns  }\SpecialCharTok{/}\NormalTok{ active\_months,}
         \AttributeTok{spend\_pm =}\NormalTok{ spend }\SpecialCharTok{/}\NormalTok{ active\_months) }\SpecialCharTok{\%\textgreater{}\%}
  \FunctionTok{left\_join}\NormalTok{(hist\_w }\SpecialCharTok{\%\textgreater{}\%} \FunctionTok{group\_by}\NormalTok{(card\_id) }\SpecialCharTok{\%\textgreater{}\%}
              \FunctionTok{summarise}\NormalTok{(}\AttributeTok{decline\_rate =} \FunctionTok{mean}\NormalTok{(}\SpecialCharTok{!}\NormalTok{approved), }\AttributeTok{.groups =} \StringTok{"drop"}\NormalTok{), }\AttributeTok{by =} \StringTok{"card\_id"}\NormalTok{)}

\NormalTok{card }\SpecialCharTok{\%\textgreater{}\%} \FunctionTok{summarise}\NormalTok{(}
  \AttributeTok{Cards =} \FunctionTok{n}\NormalTok{(),}
  \StringTok{\textasciigrave{}}\AttributeTok{Median transactions}\StringTok{\textasciigrave{}} \OtherTok{=} \FunctionTok{median}\NormalTok{(txns),}
  \StringTok{\textasciigrave{}}\AttributeTok{Median spend (BRL)}\StringTok{\textasciigrave{}} \OtherTok{=} \FunctionTok{round}\NormalTok{(}\FunctionTok{median}\NormalTok{(spend)),}
  \StringTok{\textasciigrave{}}\AttributeTok{Median ticket (BRL)}\StringTok{\textasciigrave{}} \OtherTok{=} \FunctionTok{round}\NormalTok{(}\FunctionTok{median}\NormalTok{(avg\_ticket), }\DecValTok{2}\NormalTok{),}
  \StringTok{\textasciigrave{}}\AttributeTok{Median distinct merchants}\StringTok{\textasciigrave{}} \OtherTok{=} \FunctionTok{median}\NormalTok{(merchants),}
  \StringTok{\textasciigrave{}}\AttributeTok{Median decline rate}\StringTok{\textasciigrave{}} \OtherTok{=} \FunctionTok{percent}\NormalTok{(}\FunctionTok{median}\NormalTok{(decline\_rate), }\FloatTok{0.1}\NormalTok{)) }\SpecialCharTok{\%\textgreater{}\%}
  \FunctionTok{t}\NormalTok{() }\SpecialCharTok{\%\textgreater{}\%} \FunctionTok{kable}\NormalTok{(}\AttributeTok{col.names =} \StringTok{"Value"}\NormalTok{)}
\end{Highlighting}
\end{Shaded}

{\def\LTcaptype{none} % do not increment counter
\begin{longtable}[]{@{}ll@{}}
\toprule\noalign{}
& Value \\
\midrule\noalign{}
\endhead
\bottomrule\noalign{}
\endlastfoot
Cards & 463 \\
Median transactions & 148 \\
Median spend (BRL) & 10020 \\
Median ticket (BRL) & 67.38 \\
Median distinct merchants & 56 \\
Median decline rate & 6.3\% \\
\end{longtable}
}

\subsection{The portfolio is a Pareto, and a steep
one}\label{the-portfolio-is-a-pareto-and-a-steep-one}

\begin{Shaded}
\begin{Highlighting}[]
\NormalTok{card\_sorted }\OtherTok{\textless{}{-}}\NormalTok{ card }\SpecialCharTok{\%\textgreater{}\%} \FunctionTok{arrange}\NormalTok{(}\FunctionTok{desc}\NormalTok{(spend)) }\SpecialCharTok{\%\textgreater{}\%}
  \FunctionTok{mutate}\NormalTok{(}\AttributeTok{rank =} \FunctionTok{row\_number}\NormalTok{(), }\AttributeTok{cum\_share =} \FunctionTok{cumsum}\NormalTok{(spend) }\SpecialCharTok{/} \FunctionTok{sum}\NormalTok{(spend),}
         \AttributeTok{card\_share =}\NormalTok{ rank }\SpecialCharTok{/} \FunctionTok{n}\NormalTok{())}

\FunctionTok{data.frame}\NormalTok{(}\StringTok{\textasciigrave{}}\AttributeTok{Top share of cards}\StringTok{\textasciigrave{}} \OtherTok{=} \FunctionTok{percent}\NormalTok{(}\FunctionTok{c}\NormalTok{(.}\DecValTok{05}\NormalTok{, .}\DecValTok{10}\NormalTok{, .}\DecValTok{20}\NormalTok{, .}\DecValTok{50}\NormalTok{)),}
           \StringTok{\textasciigrave{}}\AttributeTok{Share of total spend}\StringTok{\textasciigrave{}} \OtherTok{=} \FunctionTok{percent}\NormalTok{(}
             \FunctionTok{sapply}\NormalTok{(}\FunctionTok{c}\NormalTok{(.}\DecValTok{05}\NormalTok{, .}\DecValTok{10}\NormalTok{, .}\DecValTok{20}\NormalTok{, .}\DecValTok{50}\NormalTok{),}
                    \ControlFlowTok{function}\NormalTok{(p) card\_sorted}\SpecialCharTok{$}\NormalTok{cum\_share[}\FunctionTok{ceiling}\NormalTok{(p }\SpecialCharTok{*} \FunctionTok{nrow}\NormalTok{(card\_sorted))]), }\FloatTok{0.1}\NormalTok{),}
           \AttributeTok{check.names =} \ConstantTok{FALSE}\NormalTok{) }\SpecialCharTok{\%\textgreater{}\%} \FunctionTok{kable}\NormalTok{()}
\end{Highlighting}
\end{Shaded}

{\def\LTcaptype{none} % do not increment counter
\begin{longtable}[]{@{}ll@{}}
\toprule\noalign{}
Top share of cards & Share of total spend \\
\midrule\noalign{}
\endhead
\bottomrule\noalign{}
\endlastfoot
5\% & 23.5\% \\
10\% & 37.5\% \\
20\% & 57.7\% \\
50\% & 86.2\% \\
\end{longtable}
}

\begin{Shaded}
\begin{Highlighting}[]
\FunctionTok{ggplot}\NormalTok{(card\_sorted, }\FunctionTok{aes}\NormalTok{(card\_share, cum\_share)) }\SpecialCharTok{+}
  \FunctionTok{geom\_area}\NormalTok{(}\AttributeTok{fill =}\NormalTok{ PAL[}\DecValTok{3}\NormalTok{], }\AttributeTok{alpha =}\NormalTok{ .}\DecValTok{5}\NormalTok{) }\SpecialCharTok{+}
  \FunctionTok{geom\_line}\NormalTok{(}\AttributeTok{colour =}\NormalTok{ PAL[}\DecValTok{1}\NormalTok{], }\AttributeTok{size =} \DecValTok{1}\NormalTok{) }\SpecialCharTok{+}
  \FunctionTok{geom\_abline}\NormalTok{(}\AttributeTok{linetype =} \StringTok{"dashed"}\NormalTok{, }\AttributeTok{colour =} \StringTok{"grey60"}\NormalTok{) }\SpecialCharTok{+}
  \FunctionTok{geom\_vline}\NormalTok{(}\AttributeTok{xintercept =}\NormalTok{ .}\DecValTok{2}\NormalTok{, }\AttributeTok{linetype =} \StringTok{"dotted"}\NormalTok{, }\AttributeTok{colour =}\NormalTok{ PAL[}\DecValTok{5}\NormalTok{]) }\SpecialCharTok{+}
  \FunctionTok{annotate}\NormalTok{(}\StringTok{"text"}\NormalTok{, }\AttributeTok{x =}\NormalTok{ .}\DecValTok{24}\NormalTok{, }\AttributeTok{y =}\NormalTok{ .}\DecValTok{35}\NormalTok{, }\AttributeTok{hjust =} \DecValTok{0}\NormalTok{, }\AttributeTok{size =} \FloatTok{3.4}\NormalTok{, }\AttributeTok{colour =}\NormalTok{ PAL[}\DecValTok{5}\NormalTok{],}
           \AttributeTok{label =} \StringTok{"Top 20\% of cards}\SpecialCharTok{\textbackslash{}n}\StringTok{carry 58\% of spend"}\NormalTok{) }\SpecialCharTok{+}
  \FunctionTok{scale\_x\_continuous}\NormalTok{(}\AttributeTok{labels =}\NormalTok{ percent) }\SpecialCharTok{+} \FunctionTok{scale\_y\_continuous}\NormalTok{(}\AttributeTok{labels =}\NormalTok{ percent) }\SpecialCharTok{+}
  \FunctionTok{labs}\NormalTok{(}\AttributeTok{title =} \StringTok{"Spend is owned by a fifth of the portfolio"}\NormalTok{,}
       \AttributeTok{subtitle =} \StringTok{"Cumulative share of card spend, cards ranked highest to lowest"}\NormalTok{,}
       \AttributeTok{x =} \StringTok{"Share of cards"}\NormalTok{, }\AttributeTok{y =} \StringTok{"Share of spend"}\NormalTok{,}
       \AttributeTok{caption =} \StringTok{"Historical ledger, approved transactions, top 0.1\% of tickets excluded"}\NormalTok{)}
\end{Highlighting}
\end{Shaded}

\begin{center}\includegraphics{Elo_Analysis_files/figure-latex/pareto-1} \end{center}

The top 20\% of cards carry \textbf{58\%} of spend; the bottom half
carry \textbf{14\%}. The commercial reading is blunt: a promotions
programme run as a broadcast is spending most of its budget on
cardholders who contribute almost nothing, and is under-serving the
group that funds the network.

\subsection{Four cardholders, not one}\label{four-cardholders-not-one}

We segment on the two dimensions that a promotions team can actually act
on --- \textbf{value per active month} and \textbf{merchant breadth}
(how many different merchants a card touches). Both are split at the
portfolio median, which keeps the rule transparent and reproducible.

\begin{Shaded}
\begin{Highlighting}[]
\NormalTok{card }\OtherTok{\textless{}{-}}\NormalTok{ card }\SpecialCharTok{\%\textgreater{}\%}
  \FunctionTok{mutate}\NormalTok{(}\AttributeTok{segment =} \FunctionTok{case\_when}\NormalTok{(}
\NormalTok{    spend\_pm }\SpecialCharTok{\textgreater{}=} \FunctionTok{median}\NormalTok{(spend\_pm) }\SpecialCharTok{\&}\NormalTok{ merchants }\SpecialCharTok{\textgreater{}=} \FunctionTok{median}\NormalTok{(merchants) }\SpecialCharTok{\textasciitilde{}} \StringTok{"Core Everyday"}\NormalTok{,}
\NormalTok{    spend\_pm }\SpecialCharTok{\textgreater{}=} \FunctionTok{median}\NormalTok{(spend\_pm)                                   }\SpecialCharTok{\textasciitilde{}} \StringTok{"Big{-}Ticket Narrow"}\NormalTok{,}
\NormalTok{    merchants }\SpecialCharTok{\textgreater{}=} \FunctionTok{median}\NormalTok{(merchants)                                 }\SpecialCharTok{\textasciitilde{}} \StringTok{"Explorer Light"}\NormalTok{,}
    \ConstantTok{TRUE}                                                           \SpecialCharTok{\textasciitilde{}} \StringTok{"Occasional"}\NormalTok{),}
    \AttributeTok{segment =} \FunctionTok{factor}\NormalTok{(segment, }\AttributeTok{levels =} \FunctionTok{c}\NormalTok{(}\StringTok{"Core Everyday"}\NormalTok{, }\StringTok{"Big{-}Ticket Narrow"}\NormalTok{,}
                                         \StringTok{"Explorer Light"}\NormalTok{, }\StringTok{"Occasional"}\NormalTok{)))}

\NormalTok{seg\_profile }\OtherTok{\textless{}{-}}\NormalTok{ card }\SpecialCharTok{\%\textgreater{}\%} \FunctionTok{group\_by}\NormalTok{(segment) }\SpecialCharTok{\%\textgreater{}\%}
  \FunctionTok{summarise}\NormalTok{(}\AttributeTok{Cards =} \FunctionTok{n}\NormalTok{(),}
            \StringTok{\textasciigrave{}}\AttributeTok{\% of cards}\StringTok{\textasciigrave{}} \OtherTok{=} \FunctionTok{n}\NormalTok{() }\SpecialCharTok{/} \FunctionTok{nrow}\NormalTok{(card),}
            \StringTok{\textasciigrave{}}\AttributeTok{\% of spend}\StringTok{\textasciigrave{}} \OtherTok{=} \FunctionTok{sum}\NormalTok{(spend) }\SpecialCharTok{/} \FunctionTok{sum}\NormalTok{(card}\SpecialCharTok{$}\NormalTok{spend),}
            \StringTok{\textasciigrave{}}\AttributeTok{Spend / month}\StringTok{\textasciigrave{}} \OtherTok{=} \FunctionTok{mean}\NormalTok{(spend\_pm),}
            \StringTok{\textasciigrave{}}\AttributeTok{Txns / month}\StringTok{\textasciigrave{}} \OtherTok{=} \FunctionTok{mean}\NormalTok{(txns\_pm),}
            \StringTok{\textasciigrave{}}\AttributeTok{Avg ticket}\StringTok{\textasciigrave{}} \OtherTok{=} \FunctionTok{mean}\NormalTok{(avg\_ticket),}
            \StringTok{\textasciigrave{}}\AttributeTok{Merchants}\StringTok{\textasciigrave{}} \OtherTok{=} \FunctionTok{mean}\NormalTok{(merchants),}
            \StringTok{\textasciigrave{}}\AttributeTok{Days since last txn}\StringTok{\textasciigrave{}} \OtherTok{=} \FunctionTok{mean}\NormalTok{(recency\_days),}
            \StringTok{\textasciigrave{}}\AttributeTok{Decline rate}\StringTok{\textasciigrave{}} \OtherTok{=} \FunctionTok{mean}\NormalTok{(decline\_rate), }\AttributeTok{.groups =} \StringTok{"drop"}\NormalTok{)}

\NormalTok{seg\_profile }\SpecialCharTok{\%\textgreater{}\%}
  \FunctionTok{mutate}\NormalTok{(}\StringTok{\textasciigrave{}}\AttributeTok{\% of cards}\StringTok{\textasciigrave{}} \OtherTok{=} \FunctionTok{percent}\NormalTok{(}\StringTok{\textasciigrave{}}\AttributeTok{\% of cards}\StringTok{\textasciigrave{}}\NormalTok{, }\FloatTok{0.1}\NormalTok{),}
         \StringTok{\textasciigrave{}}\AttributeTok{\% of spend}\StringTok{\textasciigrave{}} \OtherTok{=} \FunctionTok{percent}\NormalTok{(}\StringTok{\textasciigrave{}}\AttributeTok{\% of spend}\StringTok{\textasciigrave{}}\NormalTok{, }\FloatTok{0.1}\NormalTok{),}
         \StringTok{\textasciigrave{}}\AttributeTok{Decline rate}\StringTok{\textasciigrave{}} \OtherTok{=} \FunctionTok{percent}\NormalTok{(}\StringTok{\textasciigrave{}}\AttributeTok{Decline rate}\StringTok{\textasciigrave{}}\NormalTok{, }\FloatTok{0.1}\NormalTok{),}
         \FunctionTok{across}\NormalTok{(}\FunctionTok{where}\NormalTok{(is.numeric), }\SpecialCharTok{\textasciitilde{}}\FunctionTok{round}\NormalTok{(.x, }\DecValTok{1}\NormalTok{))) }\SpecialCharTok{\%\textgreater{}\%} \FunctionTok{kable}\NormalTok{()}
\end{Highlighting}
\end{Shaded}

{\def\LTcaptype{none} % do not increment counter
\begin{longtable}[]{@{}
  >{\raggedright\arraybackslash}p{(\linewidth - 18\tabcolsep) * \real{0.1417}}
  >{\raggedleft\arraybackslash}p{(\linewidth - 18\tabcolsep) * \real{0.0472}}
  >{\raggedright\arraybackslash}p{(\linewidth - 18\tabcolsep) * \real{0.0866}}
  >{\raggedright\arraybackslash}p{(\linewidth - 18\tabcolsep) * \real{0.0866}}
  >{\raggedleft\arraybackslash}p{(\linewidth - 18\tabcolsep) * \real{0.1102}}
  >{\raggedleft\arraybackslash}p{(\linewidth - 18\tabcolsep) * \real{0.1024}}
  >{\raggedleft\arraybackslash}p{(\linewidth - 18\tabcolsep) * \real{0.0866}}
  >{\raggedleft\arraybackslash}p{(\linewidth - 18\tabcolsep) * \real{0.0787}}
  >{\raggedleft\arraybackslash}p{(\linewidth - 18\tabcolsep) * \real{0.1575}}
  >{\raggedright\arraybackslash}p{(\linewidth - 18\tabcolsep) * \real{0.1024}}@{}}
\toprule\noalign{}
\begin{minipage}[b]{\linewidth}\raggedright
segment
\end{minipage} & \begin{minipage}[b]{\linewidth}\raggedleft
Cards
\end{minipage} & \begin{minipage}[b]{\linewidth}\raggedright
\% of cards
\end{minipage} & \begin{minipage}[b]{\linewidth}\raggedright
\% of spend
\end{minipage} & \begin{minipage}[b]{\linewidth}\raggedleft
Spend / month
\end{minipage} & \begin{minipage}[b]{\linewidth}\raggedleft
Txns / month
\end{minipage} & \begin{minipage}[b]{\linewidth}\raggedleft
Avg ticket
\end{minipage} & \begin{minipage}[b]{\linewidth}\raggedleft
Merchants
\end{minipage} & \begin{minipage}[b]{\linewidth}\raggedleft
Days since last txn
\end{minipage} & \begin{minipage}[b]{\linewidth}\raggedright
Decline rate
\end{minipage} \\
\midrule\noalign{}
\endhead
\bottomrule\noalign{}
\endlastfoot
Core Everyday & 171 & 36.9\% & 72.3\% & 3583.9 & 35.2 & 111.8 & 114.9 &
22.2 & 5.6\% \\
Big-Ticket Narrow & 61 & 13.2\% & 12.5\% & 2459.4 & 20.3 & 148.2 & 39.9
& 44.5 & 8.3\% \\
Explorer Light & 66 & 14.3\% & 6.6\% & 655.8 & 17.4 & 40.7 & 73.6 & 24.6
& 8.5\% \\
Occasional & 165 & 35.6\% & 8.6\% & 457.9 & 9.4 & 68.2 & 28.5 & 37.9 &
12.6\% \\
\end{longtable}
}

\begin{Shaded}
\begin{Highlighting}[]
\NormalTok{seg\_profile }\SpecialCharTok{\%\textgreater{}\%} \FunctionTok{select}\NormalTok{(segment, }\StringTok{\textasciigrave{}}\AttributeTok{\% of cards}\StringTok{\textasciigrave{}}\NormalTok{, }\StringTok{\textasciigrave{}}\AttributeTok{\% of spend}\StringTok{\textasciigrave{}}\NormalTok{) }\SpecialCharTok{\%\textgreater{}\%}
  \FunctionTok{pivot\_longer}\NormalTok{(}\SpecialCharTok{{-}}\NormalTok{segment) }\SpecialCharTok{\%\textgreater{}\%}
  \FunctionTok{ggplot}\NormalTok{(}\FunctionTok{aes}\NormalTok{(segment, value, }\AttributeTok{fill =}\NormalTok{ name)) }\SpecialCharTok{+}
  \FunctionTok{geom\_col}\NormalTok{(}\AttributeTok{position =} \StringTok{"dodge"}\NormalTok{) }\SpecialCharTok{+}
  \FunctionTok{geom\_text}\NormalTok{(}\FunctionTok{aes}\NormalTok{(}\AttributeTok{label =} \FunctionTok{percent}\NormalTok{(value, }\FloatTok{0.1}\NormalTok{)), }\AttributeTok{position =} \FunctionTok{position\_dodge}\NormalTok{(.}\DecValTok{9}\NormalTok{),}
            \AttributeTok{vjust =} \SpecialCharTok{{-}}\FloatTok{0.4}\NormalTok{, }\AttributeTok{size =} \FloatTok{3.2}\NormalTok{) }\SpecialCharTok{+}
  \FunctionTok{scale\_y\_continuous}\NormalTok{(}\AttributeTok{labels =}\NormalTok{ percent, }\AttributeTok{expand =} \FunctionTok{expansion}\NormalTok{(}\AttributeTok{mult =} \FunctionTok{c}\NormalTok{(}\DecValTok{0}\NormalTok{, .}\DecValTok{15}\NormalTok{))) }\SpecialCharTok{+}
  \FunctionTok{scale\_fill\_manual}\NormalTok{(}\AttributeTok{values =} \FunctionTok{c}\NormalTok{(PAL[}\DecValTok{3}\NormalTok{], PAL[}\DecValTok{1}\NormalTok{]), }\AttributeTok{name =} \ConstantTok{NULL}\NormalTok{) }\SpecialCharTok{+}
  \FunctionTok{labs}\NormalTok{(}\AttributeTok{title =} \StringTok{"One segment is the network"}\NormalTok{,}
       \AttributeTok{subtitle =} \StringTok{"Core Everyday cardholders are a third of the base and roughly three quarters of the money"}\NormalTok{,}
       \AttributeTok{x =} \ConstantTok{NULL}\NormalTok{, }\AttributeTok{y =} \ConstantTok{NULL}\NormalTok{)}
\end{Highlighting}
\end{Shaded}

\begin{center}\includegraphics{Elo_Analysis_files/figure-latex/segments-1} \end{center}

Read across the table rather than down it:

\begin{itemize}
\tightlist
\item
  \textbf{Core Everyday} --- a third of cards, about \textbf{72\% of
  spend}, 35 transactions a month across \textasciitilde115 different
  merchants, and the \emph{lowest} decline rate. This is the card in
  someone's pocket every day.
\item
  \textbf{Big-Ticket Narrow} --- high value, but concentrated in
  \textasciitilde40 merchants and last seen \textasciitilde43 days ago.
  High revenue, thin relationship: the segment most exposed to a
  competitor's offer.
\item
  \textbf{Occasional} --- a third of cards, under 9\% of spend, and the
  \textbf{highest decline rate} of the four. A broadcast promotion
  spends a third of its budget here.
\item
  \textbf{Explorer Light} --- broad but cheap: a R\$41 ticket across 74
  merchants. Curious, low-value, and the natural audience for discovery
  offers.
\end{itemize}

\begin{Shaded}
\begin{Highlighting}[]
\NormalTok{km }\OtherTok{\textless{}{-}} \FunctionTok{kmeans}\NormalTok{(}\FunctionTok{scale}\NormalTok{(card }\SpecialCharTok{\%\textgreater{}\%} \FunctionTok{transmute}\NormalTok{(}\FunctionTok{log1p}\NormalTok{(spend\_pm), }\FunctionTok{log1p}\NormalTok{(txns\_pm),}
                                      \FunctionTok{log1p}\NormalTok{(merchants), }\FunctionTok{log1p}\NormalTok{(avg\_ticket))),}
             \AttributeTok{centers =} \DecValTok{4}\NormalTok{, }\AttributeTok{nstart =} \DecValTok{25}\NormalTok{)}
\FunctionTok{cat}\NormalTok{(}\StringTok{"Unsupervised cross{-}check — k{-}means (k = 4) cluster sizes:"}\NormalTok{,}
    \FunctionTok{paste}\NormalTok{(km}\SpecialCharTok{$}\NormalTok{size, }\AttributeTok{collapse =} \StringTok{" / "}\NormalTok{),}
    \StringTok{"}\SpecialCharTok{\textbackslash{}n}\StringTok{Between{-}cluster variance explained:"}\NormalTok{, }\FunctionTok{percent}\NormalTok{(km}\SpecialCharTok{$}\NormalTok{betweenss }\SpecialCharTok{/}\NormalTok{ km}\SpecialCharTok{$}\NormalTok{totss, }\FloatTok{0.1}\NormalTok{))}
\end{Highlighting}
\end{Shaded}

\begin{verbatim}
## Unsupervised cross-check — k-means (k = 4) cluster sizes: 138 / 105 / 149 / 71 
## Between-cluster variance explained: 65.5%
\end{verbatim}

An unsupervised k-means on the same behavioural variables separates the
base into groups of comparable size and explains a large share of the
variance, so the four-way split is a property of the data rather than an
artefact of where we drew the median.

\begin{center}\rule{0.5\linewidth}{0.5pt}\end{center}

\section{How is the network actually
used?}\label{how-is-the-network-actually-used}

\subsection{This is the identity
answer}\label{this-is-the-identity-answer}

\begin{Shaded}
\begin{Highlighting}[]
\NormalTok{instal\_mix }\OtherTok{\textless{}{-}}\NormalTok{ hist\_w }\SpecialCharTok{\%\textgreater{}\%} \FunctionTok{filter}\NormalTok{(}\SpecialCharTok{!}\FunctionTok{is.na}\NormalTok{(instal\_band)) }\SpecialCharTok{\%\textgreater{}\%}
  \FunctionTok{count}\NormalTok{(instal\_band) }\SpecialCharTok{\%\textgreater{}\%} \FunctionTok{mutate}\NormalTok{(}\AttributeTok{share =}\NormalTok{ n }\SpecialCharTok{/} \FunctionTok{sum}\NormalTok{(n))}

\NormalTok{instal\_mix }\SpecialCharTok{\%\textgreater{}\%} \FunctionTok{mutate}\NormalTok{(}\AttributeTok{share =} \FunctionTok{percent}\NormalTok{(share, }\FloatTok{0.1}\NormalTok{)) }\SpecialCharTok{\%\textgreater{}\%}
  \FunctionTok{kable}\NormalTok{(}\AttributeTok{col.names =} \FunctionTok{c}\NormalTok{(}\StringTok{"Instalments"}\NormalTok{, }\StringTok{"Transactions"}\NormalTok{, }\StringTok{"Share"}\NormalTok{),}
        \AttributeTok{format.args =} \FunctionTok{list}\NormalTok{(}\AttributeTok{big.mark =} \StringTok{","}\NormalTok{))}
\end{Highlighting}
\end{Shaded}

{\def\LTcaptype{none} % do not increment counter
\begin{longtable}[]{@{}lrl@{}}
\toprule\noalign{}
Instalments & Transactions & Share \\
\midrule\noalign{}
\endhead
\bottomrule\noalign{}
\endlastfoot
0 (lump sum) & 55,991 & 56.2\% \\
1 (single) & 39,736 & 39.9\% \\
2-3 & 2,541 & 2.6\% \\
4-6 & 914 & 0.9\% \\
7-12 & 389 & 0.4\% \\
\end{longtable}
}

\begin{Shaded}
\begin{Highlighting}[]
\FunctionTok{ggplot}\NormalTok{(instal\_mix, }\FunctionTok{aes}\NormalTok{(instal\_band, share)) }\SpecialCharTok{+}
  \FunctionTok{geom\_col}\NormalTok{(}\AttributeTok{fill =}\NormalTok{ PAL[}\DecValTok{1}\NormalTok{]) }\SpecialCharTok{+}
  \FunctionTok{geom\_text}\NormalTok{(}\FunctionTok{aes}\NormalTok{(}\AttributeTok{label =} \FunctionTok{percent}\NormalTok{(share, }\FloatTok{0.1}\NormalTok{)), }\AttributeTok{vjust =} \SpecialCharTok{{-}}\FloatTok{0.4}\NormalTok{, }\AttributeTok{size =} \FloatTok{3.6}\NormalTok{) }\SpecialCharTok{+}
  \FunctionTok{scale\_y\_continuous}\NormalTok{(}\AttributeTok{labels =}\NormalTok{ percent, }\AttributeTok{expand =} \FunctionTok{expansion}\NormalTok{(}\AttributeTok{mult =} \FunctionTok{c}\NormalTok{(}\DecValTok{0}\NormalTok{, .}\DecValTok{15}\NormalTok{))) }\SpecialCharTok{+}
  \FunctionTok{labs}\NormalTok{(}\AttributeTok{title =} \StringTok{"Elo is not an instalment network"}\NormalTok{,}
       \AttributeTok{subtitle =} \StringTok{"96\% of transactions settle in a single payment {-} the parcelado market is happening somewhere else"}\NormalTok{,}
       \AttributeTok{x =} \StringTok{"Instalments on the transaction"}\NormalTok{, }\AttributeTok{y =} \StringTok{"Share of transactions"}\NormalTok{)}
\end{Highlighting}
\end{Shaded}

\begin{center}\includegraphics{Elo_Analysis_files/figure-latex/instalments-1} \end{center}

In Brazil, \emph{parcelado} --- splitting a purchase across monthly
instalments --- is the defining behaviour of the credit card market. On
this ledger it is almost absent. \textbf{56\%} of transactions are lump
sum, \textbf{40\%} are a single payment, and everything from two
instalments upward accounts for under \textbf{4\%}.

Combined with a median ticket of \textbf{R\$35} and \textbf{63\%} of
transactions under R\$50, the identity is settled:

\begin{quote}
\textbf{Elo is a high-frequency, small-ticket, everyday-spend network.}
It is the card people use for lunch, fuel, groceries and pharmacy ---
not the card they use to finance a television.
\end{quote}

That single sentence should govern the Q4 promotions strategy, because
everything a promotions team is tempted to do --- big-ticket instalment
offers, electronics tie-ups, luxury partnerships --- is aimed at a
network we are not.

\begin{Shaded}
\begin{Highlighting}[]
\NormalTok{hist\_ok }\SpecialCharTok{\%\textgreater{}\%} \FunctionTok{count}\NormalTok{(hour) }\SpecialCharTok{\%\textgreater{}\%} \FunctionTok{mutate}\NormalTok{(}\AttributeTok{share =}\NormalTok{ n }\SpecialCharTok{/} \FunctionTok{sum}\NormalTok{(n)) }\SpecialCharTok{\%\textgreater{}\%}
  \FunctionTok{ggplot}\NormalTok{(}\FunctionTok{aes}\NormalTok{(hour, share)) }\SpecialCharTok{+}
  \FunctionTok{geom\_col}\NormalTok{(}\AttributeTok{fill =}\NormalTok{ PAL[}\DecValTok{2}\NormalTok{]) }\SpecialCharTok{+}
  \FunctionTok{scale\_y\_continuous}\NormalTok{(}\AttributeTok{labels =}\NormalTok{ percent) }\SpecialCharTok{+} \FunctionTok{scale\_x\_continuous}\NormalTok{(}\AttributeTok{breaks =} \FunctionTok{seq}\NormalTok{(}\DecValTok{0}\NormalTok{, }\DecValTok{23}\NormalTok{, }\DecValTok{3}\NormalTok{)) }\SpecialCharTok{+}
  \FunctionTok{labs}\NormalTok{(}\AttributeTok{title =} \StringTok{"A daytime, everyday rhythm"}\NormalTok{,}
       \AttributeTok{subtitle =} \StringTok{"Share of approved transactions by hour {-} a lunchtime peak, then a long afternoon plateau"}\NormalTok{,}
       \AttributeTok{x =} \StringTok{"Hour of day"}\NormalTok{, }\AttributeTok{y =} \ConstantTok{NULL}\NormalTok{)}
\end{Highlighting}
\end{Shaded}

\begin{center}\includegraphics{Elo_Analysis_files/figure-latex/time-patterns-1} \end{center}

The clock agrees with the ticket. Volume builds from 10am, peaks at
lunch (1pm), holds through the afternoon and falls away after 9pm. This
is consumption, not commerce.

\subsection{Authorisation: a small leak in a specific
place}\label{authorisation-a-small-leak-in-a-specific-place}

\begin{Shaded}
\begin{Highlighting}[]
\NormalTok{approval\_overall }\OtherTok{\textless{}{-}} \FunctionTok{mean}\NormalTok{(hist\_txn}\SpecialCharTok{$}\NormalTok{approved)}

\NormalTok{ticket\_band }\OtherTok{\textless{}{-}}\NormalTok{ hist\_w }\SpecialCharTok{\%\textgreater{}\%}
  \FunctionTok{mutate}\NormalTok{(}\AttributeTok{band =} \FunctionTok{cut}\NormalTok{(amount\_brl, }\FunctionTok{c}\NormalTok{(}\DecValTok{0}\NormalTok{, }\DecValTok{25}\NormalTok{, }\DecValTok{50}\NormalTok{, }\DecValTok{100}\NormalTok{, }\DecValTok{250}\NormalTok{, }\DecValTok{500}\NormalTok{, }\DecValTok{1000}\NormalTok{, }\ConstantTok{Inf}\NormalTok{),}
                    \AttributeTok{labels =} \FunctionTok{c}\NormalTok{(}\StringTok{"\textless{}25"}\NormalTok{, }\StringTok{"25{-}50"}\NormalTok{, }\StringTok{"50{-}100"}\NormalTok{, }\StringTok{"100{-}250"}\NormalTok{, }\StringTok{"250{-}500"}\NormalTok{, }\StringTok{"500{-}1k"}\NormalTok{, }\StringTok{"1k+"}\NormalTok{))) }\SpecialCharTok{\%\textgreater{}\%}
  \FunctionTok{group\_by}\NormalTok{(band) }\SpecialCharTok{\%\textgreater{}\%}
  \FunctionTok{summarise}\NormalTok{(}\AttributeTok{Transactions =} \FunctionTok{n}\NormalTok{(), }\StringTok{\textasciigrave{}}\AttributeTok{\% of volume}\StringTok{\textasciigrave{}} \OtherTok{=} \FunctionTok{n}\NormalTok{() }\SpecialCharTok{/} \FunctionTok{nrow}\NormalTok{(hist\_w),}
            \StringTok{\textasciigrave{}}\AttributeTok{Decline rate}\StringTok{\textasciigrave{}} \OtherTok{=} \FunctionTok{mean}\NormalTok{(}\SpecialCharTok{!}\NormalTok{approved), }\AttributeTok{.groups =} \StringTok{"drop"}\NormalTok{)}

\FunctionTok{cat}\NormalTok{(}\StringTok{"Overall approval rate:"}\NormalTok{, }\FunctionTok{percent}\NormalTok{(approval\_overall, }\FloatTok{0.01}\NormalTok{))}
\end{Highlighting}
\end{Shaded}

\begin{verbatim}
## Overall approval rate: 93.00%
\end{verbatim}

\begin{Shaded}
\begin{Highlighting}[]
\NormalTok{ticket\_band }\SpecialCharTok{\%\textgreater{}\%} \FunctionTok{mutate}\NormalTok{(}\StringTok{\textasciigrave{}}\AttributeTok{\% of volume}\StringTok{\textasciigrave{}} \OtherTok{=} \FunctionTok{percent}\NormalTok{(}\StringTok{\textasciigrave{}}\AttributeTok{\% of volume}\StringTok{\textasciigrave{}}\NormalTok{, }\FloatTok{0.1}\NormalTok{),}
                       \StringTok{\textasciigrave{}}\AttributeTok{Decline rate}\StringTok{\textasciigrave{}} \OtherTok{=} \FunctionTok{percent}\NormalTok{(}\StringTok{\textasciigrave{}}\AttributeTok{Decline rate}\StringTok{\textasciigrave{}}\NormalTok{, }\FloatTok{0.1}\NormalTok{)) }\SpecialCharTok{\%\textgreater{}\%}
  \FunctionTok{kable}\NormalTok{(}\AttributeTok{format.args =} \FunctionTok{list}\NormalTok{(}\AttributeTok{big.mark =} \StringTok{","}\NormalTok{))}
\end{Highlighting}
\end{Shaded}

{\def\LTcaptype{none} % do not increment counter
\begin{longtable}[]{@{}lrll@{}}
\toprule\noalign{}
band & Transactions & \% of volume & Decline rate \\
\midrule\noalign{}
\endhead
\bottomrule\noalign{}
\endlastfoot
\textless25 & 39,374 & 39.4\% & 6.7\% \\
25-50 & 23,223 & 23.2\% & 6.4\% \\
50-100 & 17,506 & 17.5\% & 6.3\% \\
100-250 & 13,529 & 13.5\% & 7.2\% \\
250-500 & 3,790 & 3.8\% & 9.4\% \\
500-1k & 1,481 & 1.5\% & 13.6\% \\
1k+ & 998 & 1.0\% & 18.7\% \\
\end{longtable}
}

\begin{Shaded}
\begin{Highlighting}[]
\FunctionTok{ggplot}\NormalTok{(ticket\_band, }\FunctionTok{aes}\NormalTok{(band, }\StringTok{\textasciigrave{}}\AttributeTok{Decline rate}\StringTok{\textasciigrave{}}\NormalTok{)) }\SpecialCharTok{+}
  \FunctionTok{geom\_col}\NormalTok{(}\AttributeTok{fill =}\NormalTok{ PAL[}\DecValTok{5}\NormalTok{]) }\SpecialCharTok{+}
  \FunctionTok{geom\_text}\NormalTok{(}\FunctionTok{aes}\NormalTok{(}\AttributeTok{label =} \FunctionTok{percent}\NormalTok{(}\StringTok{\textasciigrave{}}\AttributeTok{Decline rate}\StringTok{\textasciigrave{}}\NormalTok{, }\FloatTok{0.1}\NormalTok{)), }\AttributeTok{vjust =} \SpecialCharTok{{-}}\FloatTok{0.4}\NormalTok{, }\AttributeTok{size =} \FloatTok{3.5}\NormalTok{) }\SpecialCharTok{+}
  \FunctionTok{scale\_y\_continuous}\NormalTok{(}\AttributeTok{labels =}\NormalTok{ percent, }\AttributeTok{expand =} \FunctionTok{expansion}\NormalTok{(}\AttributeTok{mult =} \FunctionTok{c}\NormalTok{(}\DecValTok{0}\NormalTok{, .}\DecValTok{18}\NormalTok{))) }\SpecialCharTok{+}
  \FunctionTok{labs}\NormalTok{(}\AttributeTok{title =} \StringTok{"Declines climb steeply with basket size"}\NormalTok{,}
       \AttributeTok{subtitle =} \StringTok{"A 6.7\% decline rate under R$25 becomes 18.7\% above R$1,000"}\NormalTok{,}
       \AttributeTok{x =} \StringTok{"Ticket size (BRL)"}\NormalTok{, }\AttributeTok{y =} \StringTok{"Declined share"}\NormalTok{)}
\end{Highlighting}
\end{Shaded}

\begin{center}\includegraphics{Elo_Analysis_files/figure-latex/decline-chart-1} \end{center}

Approval is \textbf{93.0\%} overall, which sounds healthy. It is not
evenly distributed. Below R\$25 we decline 6.7\% of attempts; above
R\$1,000 we decline \textbf{18.7\%}. The network is at its \emph{worst}
exactly where a promotion would be trying to push a cardholder to trade
up.

\subsection{What actually causes a
decline?}\label{what-actually-causes-a-decline}

\begin{Shaded}
\begin{Highlighting}[]
\NormalTok{ct\_df }\OtherTok{\textless{}{-}}\NormalTok{ hist\_w }\SpecialCharTok{\%\textgreater{}\%} \FunctionTok{filter}\NormalTok{(}\SpecialCharTok{!}\FunctionTok{is.na}\NormalTok{(instal\_band)) }\SpecialCharTok{\%\textgreater{}\%} \FunctionTok{mutate}\NormalTok{(}\AttributeTok{instal\_band =} \FunctionTok{droplevels}\NormalTok{(instal\_band))}
\NormalTok{chi }\OtherTok{\textless{}{-}} \FunctionTok{chisq.test}\NormalTok{(}\FunctionTok{table}\NormalTok{(ct\_df}\SpecialCharTok{$}\NormalTok{instal\_band, ct\_df}\SpecialCharTok{$}\NormalTok{approved))}
\FunctionTok{cat}\NormalTok{(}\StringTok{"Chi{-}square, instalment band x authorisation:  X{-}squared ="}\NormalTok{, }\FunctionTok{round}\NormalTok{(chi}\SpecialCharTok{$}\NormalTok{statistic, }\DecValTok{1}\NormalTok{),}
    \StringTok{" df ="}\NormalTok{, chi}\SpecialCharTok{$}\NormalTok{parameter, }\StringTok{" p ="}\NormalTok{, }\FunctionTok{format.pval}\NormalTok{(chi}\SpecialCharTok{$}\NormalTok{p.value, }\AttributeTok{digits =} \DecValTok{3}\NormalTok{))}
\end{Highlighting}
\end{Shaded}

\begin{verbatim}
## Chi-square, instalment band x authorisation:  X-squared = 446.1  df = 4  p = <0.0000000000000002
\end{verbatim}

\begin{Shaded}
\begin{Highlighting}[]
\NormalTok{mdl }\OtherTok{\textless{}{-}}\NormalTok{ hist\_w }\SpecialCharTok{\%\textgreater{}\%} \FunctionTok{filter}\NormalTok{(}\SpecialCharTok{!}\FunctionTok{is.na}\NormalTok{(instal)) }\SpecialCharTok{\%\textgreater{}\%}
  \FunctionTok{transmute}\NormalTok{(}\AttributeTok{decline =} \FunctionTok{as.integer}\NormalTok{(}\SpecialCharTok{!}\NormalTok{approved),}
            \AttributeTok{log\_amt =} \FunctionTok{log1p}\NormalTok{(}\FunctionTok{pmax}\NormalTok{(amount\_brl, }\FloatTok{0.01}\NormalTok{)),}
            \AttributeTok{multi   =} \FunctionTok{as.integer}\NormalTok{(instal }\SpecialCharTok{\textgreater{}=} \DecValTok{2}\NormalTok{),}
            \AttributeTok{cat1\_y  =} \FunctionTok{as.integer}\NormalTok{(category\_1 }\SpecialCharTok{==} \StringTok{"Y"}\NormalTok{))}
\NormalTok{fit }\OtherTok{\textless{}{-}} \FunctionTok{glm}\NormalTok{(decline }\SpecialCharTok{\textasciitilde{}}\NormalTok{ log\_amt }\SpecialCharTok{+}\NormalTok{ multi }\SpecialCharTok{+}\NormalTok{ cat1\_y, }\AttributeTok{data =}\NormalTok{ mdl, }\AttributeTok{family =} \FunctionTok{binomial}\NormalTok{())}

\FunctionTok{data.frame}\NormalTok{(}\AttributeTok{Term =} \FunctionTok{c}\NormalTok{(}\StringTok{"(Intercept)"}\NormalTok{, }\StringTok{"log(amount)"}\NormalTok{, }\StringTok{"2+ instalments"}\NormalTok{, }\StringTok{"category\_1 = Y"}\NormalTok{),}
           \AttributeTok{Coefficient =} \FunctionTok{round}\NormalTok{(}\FunctionTok{coef}\NormalTok{(fit), }\DecValTok{4}\NormalTok{),}
           \StringTok{\textasciigrave{}}\AttributeTok{Odds ratio}\StringTok{\textasciigrave{}} \OtherTok{=} \FunctionTok{round}\NormalTok{(}\FunctionTok{exp}\NormalTok{(}\FunctionTok{coef}\NormalTok{(fit)), }\DecValTok{3}\NormalTok{),}
           \StringTok{\textasciigrave{}}\AttributeTok{z value}\StringTok{\textasciigrave{}} \OtherTok{=} \FunctionTok{round}\NormalTok{(}\FunctionTok{summary}\NormalTok{(fit)}\SpecialCharTok{$}\NormalTok{coefficients[, }\DecValTok{3}\NormalTok{], }\DecValTok{1}\NormalTok{),}
           \AttributeTok{check.names =} \ConstantTok{FALSE}\NormalTok{, }\AttributeTok{row.names =} \ConstantTok{NULL}\NormalTok{) }\SpecialCharTok{\%\textgreater{}\%} \FunctionTok{kable}\NormalTok{()}
\end{Highlighting}
\end{Shaded}

{\def\LTcaptype{none} % do not increment counter
\begin{longtable}[]{@{}lrrr@{}}
\toprule\noalign{}
Term & Coefficient & Odds ratio & z value \\
\midrule\noalign{}
\endhead
\bottomrule\noalign{}
\endlastfoot
(Intercept) & -3.0172 & 0.049 & -69.8 \\
log(amount) & 0.0714 & 1.074 & 6.4 \\
2+ instalments & 0.1148 & 1.122 & 2.0 \\
category\_1 = Y & 1.4473 & 4.252 & 41.1 \\
\end{longtable}
}

\begin{Shaded}
\begin{Highlighting}[]
\FunctionTok{cat}\NormalTok{(}\StringTok{"McFadden pseudo R{-}squared:"}\NormalTok{, }\FunctionTok{round}\NormalTok{(}\DecValTok{1} \SpecialCharTok{{-}}\NormalTok{ fit}\SpecialCharTok{$}\NormalTok{deviance }\SpecialCharTok{/}\NormalTok{ fit}\SpecialCharTok{$}\NormalTok{null.deviance, }\DecValTok{4}\NormalTok{))}
\end{Highlighting}
\end{Shaded}

\begin{verbatim}
## McFadden pseudo R-squared: 0.0313
\end{verbatim}

The association between instalments and authorisation is real and not a
fluke of sample size. But the logistic model says something more useful:
the dominant driver is not the amount and not the instalment plan --- it
is the anonymised flag \textbf{\texttt{category\_1\ =\ Y}, which
multiplies the odds of a decline by roughly 4.3×}. Only 6.1\% of
transactions carry that flag, and they are being turned away at more
than four times the normal rate.

\textbf{Business consequence.} Before Elo spends a cent on Q4
promotions, the risk team should be asked what
\texttt{category\_1\ =\ Y} represents. If it is a merchant risk
classification, we are currently promoting into a population we then
decline. That is a promotion budget that generates a bad cardholder
experience and no revenue.

\subsection{Growth is activation, not
deepening}\label{growth-is-activation-not-deepening}

\begin{Shaded}
\begin{Highlighting}[]
\NormalTok{monthly }\OtherTok{\textless{}{-}}\NormalTok{ hist\_ok }\SpecialCharTok{\%\textgreater{}\%} \FunctionTok{group\_by}\NormalTok{(month) }\SpecialCharTok{\%\textgreater{}\%}
  \FunctionTok{summarise}\NormalTok{(}\AttributeTok{Transactions =} \FunctionTok{n}\NormalTok{(), }\AttributeTok{Spend =} \FunctionTok{sum}\NormalTok{(amount\_brl),}
            \StringTok{\textasciigrave{}}\AttributeTok{Active cards}\StringTok{\textasciigrave{}} \OtherTok{=} \FunctionTok{n\_distinct}\NormalTok{(card\_id),}
            \StringTok{\textasciigrave{}}\AttributeTok{Avg ticket}\StringTok{\textasciigrave{}} \OtherTok{=} \FunctionTok{mean}\NormalTok{(amount\_brl), }\AttributeTok{.groups =} \StringTok{"drop"}\NormalTok{) }\SpecialCharTok{\%\textgreater{}\%}
  \FunctionTok{mutate}\NormalTok{(}\StringTok{\textasciigrave{}}\AttributeTok{Txns per active card}\StringTok{\textasciigrave{}} \OtherTok{=} \FunctionTok{round}\NormalTok{(Transactions }\SpecialCharTok{/} \StringTok{\textasciigrave{}}\AttributeTok{Active cards}\StringTok{\textasciigrave{}}\NormalTok{, }\DecValTok{1}\NormalTok{))}

\NormalTok{monthly }\SpecialCharTok{\%\textgreater{}\%} \FunctionTok{mutate}\NormalTok{(}\AttributeTok{Spend =} \FunctionTok{comma}\NormalTok{(}\FunctionTok{round}\NormalTok{(Spend)), }\StringTok{\textasciigrave{}}\AttributeTok{Avg ticket}\StringTok{\textasciigrave{}} \OtherTok{=} \FunctionTok{round}\NormalTok{(}\StringTok{\textasciigrave{}}\AttributeTok{Avg ticket}\StringTok{\textasciigrave{}}\NormalTok{, }\DecValTok{2}\NormalTok{)) }\SpecialCharTok{\%\textgreater{}\%}
  \FunctionTok{kable}\NormalTok{(}\AttributeTok{format.args =} \FunctionTok{list}\NormalTok{(}\AttributeTok{big.mark =} \StringTok{","}\NormalTok{))}
\end{Highlighting}
\end{Shaded}

{\def\LTcaptype{none} % do not increment counter
\begin{longtable}[]{@{}
  >{\raggedright\arraybackslash}p{(\linewidth - 10\tabcolsep) * \real{0.1081}}
  >{\raggedleft\arraybackslash}p{(\linewidth - 10\tabcolsep) * \real{0.1757}}
  >{\raggedright\arraybackslash}p{(\linewidth - 10\tabcolsep) * \real{0.1081}}
  >{\raggedleft\arraybackslash}p{(\linewidth - 10\tabcolsep) * \real{0.1757}}
  >{\raggedleft\arraybackslash}p{(\linewidth - 10\tabcolsep) * \real{0.1486}}
  >{\raggedleft\arraybackslash}p{(\linewidth - 10\tabcolsep) * \real{0.2838}}@{}}
\toprule\noalign{}
\begin{minipage}[b]{\linewidth}\raggedright
month
\end{minipage} & \begin{minipage}[b]{\linewidth}\raggedleft
Transactions
\end{minipage} & \begin{minipage}[b]{\linewidth}\raggedright
Spend
\end{minipage} & \begin{minipage}[b]{\linewidth}\raggedleft
Active cards
\end{minipage} & \begin{minipage}[b]{\linewidth}\raggedleft
Avg ticket
\end{minipage} & \begin{minipage}[b]{\linewidth}\raggedleft
Txns per active card
\end{minipage} \\
\midrule\noalign{}
\endhead
\bottomrule\noalign{}
\endlastfoot
2017-01 & 4,386 & 357,419 & 241 & 81.49 & 18.2 \\
2017-02 & 4,415 & 327,484 & 246 & 74.18 & 17.9 \\
2017-03 & 5,058 & 419,303 & 259 & 82.90 & 19.5 \\
2017-04 & 4,712 & 367,270 & 271 & 77.94 & 17.4 \\
2017-05 & 5,154 & 384,569 & 293 & 74.62 & 17.6 \\
2017-06 & 5,404 & 409,562 & 302 & 75.79 & 17.9 \\
2017-07 & 6,989 & 555,876 & 333 & 79.54 & 21.0 \\
2017-08 & 7,562 & 603,768 & 349 & 79.84 & 21.7 \\
2017-09 & 7,784 & 636,632 & 368 & 81.79 & 21.2 \\
2017-10 & 8,476 & 726,090 & 395 & 85.66 & 21.5 \\
2017-11 & 8,571 & 773,637 & 398 & 90.26 & 21.5 \\
2017-12 & 9,826 & 964,849 & 396 & 98.19 & 24.8 \\
2018-01 & 8,079 & 769,914 & 368 & 95.30 & 22.0 \\
2018-02 & 6,514 & 574,352 & 334 & 88.17 & 19.5 \\
\end{longtable}
}

\begin{Shaded}
\begin{Highlighting}[]
\FunctionTok{ggplot}\NormalTok{(monthly, }\FunctionTok{aes}\NormalTok{(month, Transactions, }\AttributeTok{group =} \DecValTok{1}\NormalTok{)) }\SpecialCharTok{+}
  \FunctionTok{geom\_line}\NormalTok{(}\AttributeTok{colour =}\NormalTok{ PAL[}\DecValTok{1}\NormalTok{], }\AttributeTok{size =} \DecValTok{1}\NormalTok{) }\SpecialCharTok{+} \FunctionTok{geom\_point}\NormalTok{(}\AttributeTok{colour =}\NormalTok{ PAL[}\DecValTok{1}\NormalTok{]) }\SpecialCharTok{+}
  \FunctionTok{geom\_line}\NormalTok{(}\FunctionTok{aes}\NormalTok{(}\AttributeTok{y =} \StringTok{\textasciigrave{}}\AttributeTok{Active cards}\StringTok{\textasciigrave{}} \SpecialCharTok{*} \DecValTok{20}\NormalTok{), }\AttributeTok{colour =}\NormalTok{ PAL[}\DecValTok{4}\NormalTok{], }\AttributeTok{linetype =} \StringTok{"dashed"}\NormalTok{, }\AttributeTok{size =}\NormalTok{ .}\DecValTok{9}\NormalTok{) }\SpecialCharTok{+}
  \FunctionTok{scale\_y\_continuous}\NormalTok{(}\AttributeTok{labels =}\NormalTok{ comma,}
                     \AttributeTok{sec.axis =} \FunctionTok{sec\_axis}\NormalTok{(}\SpecialCharTok{\textasciitilde{}}\NormalTok{ . }\SpecialCharTok{/} \DecValTok{20}\NormalTok{, }\AttributeTok{name =} \StringTok{"Active cards (dashed)"}\NormalTok{)) }\SpecialCharTok{+}
  \FunctionTok{labs}\NormalTok{(}\AttributeTok{title =} \StringTok{"Volume grows because cards activate, not because cards spend more"}\NormalTok{,}
       \AttributeTok{subtitle =} \StringTok{"Approved transactions (solid, left) against active cards (dashed, right)"}\NormalTok{,}
       \AttributeTok{x =} \ConstantTok{NULL}\NormalTok{, }\AttributeTok{y =} \StringTok{"Transactions"}\NormalTok{) }\SpecialCharTok{+}
  \FunctionTok{theme}\NormalTok{(}\AttributeTok{axis.text.x =} \FunctionTok{element\_text}\NormalTok{(}\AttributeTok{angle =} \DecValTok{45}\NormalTok{, }\AttributeTok{hjust =} \DecValTok{1}\NormalTok{))}
\end{Highlighting}
\end{Shaded}

\begin{center}\includegraphics{Elo_Analysis_files/figure-latex/monthly-1} \end{center}

Transactions roughly double over the window, but transactions \emph{per
active card} barely move --- 18 in January 2017, 20 in February 2018,
with a December peak of 25. Average ticket rises from R\$81 to R\$98 in
December and falls back.

That is a clean diagnosis: \textbf{the network grows by adding cards,
not by getting existing cards to do more.} The promotions programme,
whose entire job is the second thing, is currently not visible in this
series.

\begin{center}\rule{0.5\linewidth}{0.5pt}\end{center}

\section{The merchant side of the
marketplace}\label{the-merchant-side-of-the-marketplace}

\subsection{A base of very small
merchants}\label{a-base-of-very-small-merchants}

\begin{Shaded}
\begin{Highlighting}[]
\NormalTok{size\_pyramid }\OtherTok{\textless{}{-}}\NormalTok{ merch }\SpecialCharTok{\%\textgreater{}\%} \FunctionTok{filter}\NormalTok{(}\SpecialCharTok{!}\FunctionTok{is.na}\NormalTok{(most\_recent\_sales\_range)) }\SpecialCharTok{\%\textgreater{}\%}
  \FunctionTok{count}\NormalTok{(most\_recent\_sales\_range) }\SpecialCharTok{\%\textgreater{}\%} \FunctionTok{mutate}\NormalTok{(}\AttributeTok{share =}\NormalTok{ n }\SpecialCharTok{/} \FunctionTok{sum}\NormalTok{(n))}

\NormalTok{size\_pyramid }\SpecialCharTok{\%\textgreater{}\%} \FunctionTok{mutate}\NormalTok{(}\AttributeTok{share =} \FunctionTok{percent}\NormalTok{(share, }\FloatTok{0.1}\NormalTok{)) }\SpecialCharTok{\%\textgreater{}\%}
  \FunctionTok{kable}\NormalTok{(}\AttributeTok{col.names =} \FunctionTok{c}\NormalTok{(}\StringTok{"Sales range (A = largest)"}\NormalTok{, }\StringTok{"Merchants"}\NormalTok{, }\StringTok{"Share of base"}\NormalTok{),}
        \AttributeTok{format.args =} \FunctionTok{list}\NormalTok{(}\AttributeTok{big.mark =} \StringTok{","}\NormalTok{))}
\end{Highlighting}
\end{Shaded}

{\def\LTcaptype{none} % do not increment counter
\begin{longtable}[]{@{}lrl@{}}
\toprule\noalign{}
Sales range (A = largest) & Merchants & Share of base \\
\midrule\noalign{}
\endhead
\bottomrule\noalign{}
\endlastfoot
A & 991 & 0.3\% \\
B & 5,030 & 1.5\% \\
C & 34,059 & 10.2\% \\
D & 117,460 & 35.1\% \\
E & 177,093 & 52.9\% \\
\end{longtable}
}

\textbf{88\%} of the merchant base sits in the two smallest revenue
bands. Only \textbf{1,005 merchants --- 0.3\% --- are in the top band.}
Elo's marketplace is a long tail of very small businesses with a handful
of large ones at the top.

\subsection{And a ledger that only touches a fraction of
it}\label{and-a-ledger-that-only-touches-a-fraction-of-it}

\begin{Shaded}
\begin{Highlighting}[]
\DocumentationTok{\#\# 0.5\% of ledger rows carry no merchant\_id. They are dropped from every}
\DocumentationTok{\#\# merchant{-}level count, otherwise "unknown" ranks as the largest merchant.}
\NormalTok{hist\_m }\OtherTok{\textless{}{-}}\NormalTok{ hist\_ok }\SpecialCharTok{\%\textgreater{}\%} \FunctionTok{filter}\NormalTok{(}\SpecialCharTok{!}\FunctionTok{is.na}\NormalTok{(merchant\_id))}
\NormalTok{new\_m  }\OtherTok{\textless{}{-}}\NormalTok{ new\_w   }\SpecialCharTok{\%\textgreater{}\%} \FunctionTok{filter}\NormalTok{(}\SpecialCharTok{!}\FunctionTok{is.na}\NormalTok{(merchant\_id))}

\NormalTok{merch\_txn }\OtherTok{\textless{}{-}}\NormalTok{ hist\_m }\SpecialCharTok{\%\textgreater{}\%} \FunctionTok{count}\NormalTok{(merchant\_id, }\AttributeTok{name =} \StringTok{"txns"}\NormalTok{) }\SpecialCharTok{\%\textgreater{}\%}
  \FunctionTok{arrange}\NormalTok{(}\FunctionTok{desc}\NormalTok{(txns)) }\SpecialCharTok{\%\textgreater{}\%} \FunctionTok{mutate}\NormalTok{(}\AttributeTok{cum\_share =} \FunctionTok{cumsum}\NormalTok{(txns) }\SpecialCharTok{/} \FunctionTok{sum}\NormalTok{(txns))}
\NormalTok{seen }\OtherTok{\textless{}{-}} \FunctionTok{union}\NormalTok{(}\FunctionTok{unique}\NormalTok{(hist\_m}\SpecialCharTok{$}\NormalTok{merchant\_id), }\FunctionTok{unique}\NormalTok{(new\_m}\SpecialCharTok{$}\NormalTok{merchant\_id))}

\FunctionTok{data.frame}\NormalTok{(}\AttributeTok{Metric =} \FunctionTok{c}\NormalTok{(}\StringTok{"Merchants in the reference file"}\NormalTok{,}
                      \StringTok{"Merchants transacted with in the historical window"}\NormalTok{,}
                      \StringTok{"Merchants transacted with in the novelty window"}\NormalTok{,}
                      \StringTok{"Merchants never seen in either window"}\NormalTok{,}
                      \StringTok{"Share of transactions from the top 100 merchants"}\NormalTok{,}
                      \StringTok{"Share of transactions from the top 1,000 merchants"}\NormalTok{),}
           \AttributeTok{Value =} \FunctionTok{c}\NormalTok{(}\FunctionTok{comma}\NormalTok{(}\FunctionTok{nrow}\NormalTok{(merch)), }\FunctionTok{comma}\NormalTok{(}\FunctionTok{nrow}\NormalTok{(merch\_txn)),}
                     \FunctionTok{comma}\NormalTok{(}\FunctionTok{n\_distinct}\NormalTok{(new\_m}\SpecialCharTok{$}\NormalTok{merchant\_id)),}
                     \FunctionTok{percent}\NormalTok{(}\DecValTok{1} \SpecialCharTok{{-}} \FunctionTok{length}\NormalTok{(seen) }\SpecialCharTok{/} \FunctionTok{nrow}\NormalTok{(merch), }\FloatTok{0.1}\NormalTok{),}
                     \FunctionTok{percent}\NormalTok{(merch\_txn}\SpecialCharTok{$}\NormalTok{cum\_share[}\DecValTok{100}\NormalTok{], }\FloatTok{0.1}\NormalTok{),}
                     \FunctionTok{percent}\NormalTok{(merch\_txn}\SpecialCharTok{$}\NormalTok{cum\_share[}\DecValTok{1000}\NormalTok{], }\FloatTok{0.1}\NormalTok{))) }\SpecialCharTok{\%\textgreater{}\%} \FunctionTok{kable}\NormalTok{()}
\end{Highlighting}
\end{Shaded}

{\def\LTcaptype{none} % do not increment counter
\begin{longtable}[]{@{}ll@{}}
\toprule\noalign{}
Metric & Value \\
\midrule\noalign{}
\endhead
\bottomrule\noalign{}
\endlastfoot
Merchants in the reference file & 334,633 \\
Merchants transacted with in the historical window & 20,549 \\
Merchants transacted with in the novelty window & 48,941 \\
Merchants never seen in either window & 82.4\% \\
Share of transactions from the top 100 merchants & 22.6\% \\
Share of transactions from the top 1,000 merchants & 50.3\% \\
\end{longtable}
}

Two facts sit awkwardly together. \textbf{82\% of the merchants we
maintain reference data for never appear in either transaction window},
while the \textbf{top 1,000 merchants carry half of all transactions}.
The merchant file is largely an archive; the live marketplace is far
smaller and far more concentrated than the file implies.

\subsection{Health of the supply base}\label{health-of-the-supply-base}

\begin{Shaded}
\begin{Highlighting}[]
\NormalTok{merch\_seg }\OtherTok{\textless{}{-}}\NormalTok{ merch }\SpecialCharTok{\%\textgreater{}\%}
  \FunctionTok{filter}\NormalTok{(}\FunctionTok{is.finite}\NormalTok{(avg\_purchases\_lag3), }\FunctionTok{is.finite}\NormalTok{(avg\_purchases\_lag12),}
\NormalTok{         avg\_purchases\_lag12 }\SpecialCharTok{!=} \DecValTok{0}\NormalTok{) }\SpecialCharTok{\%\textgreater{}\%}
  \FunctionTok{mutate}\NormalTok{(}\AttributeTok{momentum =}\NormalTok{ avg\_purchases\_lag3 }\SpecialCharTok{/}\NormalTok{ avg\_purchases\_lag12,}
         \AttributeTok{health =} \FunctionTok{case\_when}\NormalTok{(active\_months\_lag12 }\SpecialCharTok{\textless{}} \DecValTok{12} \SpecialCharTok{\textasciitilde{}} \StringTok{"Intermittent"}\NormalTok{,}
\NormalTok{                            momentum }\SpecialCharTok{\textgreater{}=} \FloatTok{1.05} \SpecialCharTok{\textasciitilde{}} \StringTok{"Accelerating"}\NormalTok{,}
\NormalTok{                            momentum }\SpecialCharTok{\textless{}=} \FloatTok{0.95} \SpecialCharTok{\textasciitilde{}} \StringTok{"Decelerating"}\NormalTok{,}
                            \ConstantTok{TRUE} \SpecialCharTok{\textasciitilde{}} \StringTok{"Steady"}\NormalTok{))}

\NormalTok{merch\_seg }\SpecialCharTok{\%\textgreater{}\%} \FunctionTok{count}\NormalTok{(health) }\SpecialCharTok{\%\textgreater{}\%} \FunctionTok{mutate}\NormalTok{(}\AttributeTok{Share =} \FunctionTok{percent}\NormalTok{(n }\SpecialCharTok{/} \FunctionTok{sum}\NormalTok{(n), }\FloatTok{0.1}\NormalTok{)) }\SpecialCharTok{\%\textgreater{}\%}
  \FunctionTok{arrange}\NormalTok{(}\FunctionTok{desc}\NormalTok{(n)) }\SpecialCharTok{\%\textgreater{}\%} \FunctionTok{kable}\NormalTok{(}\AttributeTok{col.names =} \FunctionTok{c}\NormalTok{(}\StringTok{"Health"}\NormalTok{, }\StringTok{"Merchants"}\NormalTok{, }\StringTok{"Share"}\NormalTok{),}
                             \AttributeTok{format.args =} \FunctionTok{list}\NormalTok{(}\AttributeTok{big.mark =} \StringTok{","}\NormalTok{))}
\end{Highlighting}
\end{Shaded}

{\def\LTcaptype{none} % do not increment counter
\begin{longtable}[]{@{}lrl@{}}
\toprule\noalign{}
Health & Merchants & Share \\
\midrule\noalign{}
\endhead
\bottomrule\noalign{}
\endlastfoot
Decelerating & 129,755 & 38.8\% \\
Accelerating & 91,276 & 27.3\% \\
Steady & 84,015 & 25.1\% \\
Intermittent & 29,584 & 8.8\% \\
\end{longtable}
}

Comparing each merchant's last-three-month transaction ratio with its
last-twelve-month ratio splits the base into four states. \textbf{39\%
are decelerating} and a further \textbf{9\% transact intermittently} ---
nearly half the marketplace is flat or shrinking. Only 27\% are
accelerating. A promotions programme has a genuine supply-side job here,
but only if it is aimed at merchants that can convert traffic, not at
ones already on the way out.

\subsection{Where the money actually
lands}\label{where-the-money-actually-lands}

\begin{Shaded}
\begin{Highlighting}[]
\NormalTok{tier }\OtherTok{\textless{}{-}}\NormalTok{ merch }\SpecialCharTok{\%\textgreater{}\%} \FunctionTok{select}\NormalTok{(merchant\_id, most\_recent\_sales\_range)}
\NormalTok{tier\_flow }\OtherTok{\textless{}{-}} \FunctionTok{bind\_rows}\NormalTok{(}
\NormalTok{  hist\_m }\SpecialCharTok{\%\textgreater{}\%} \FunctionTok{select}\NormalTok{(merchant\_id) }\SpecialCharTok{\%\textgreater{}\%} \FunctionTok{mutate}\NormalTok{(}\AttributeTok{window =} \StringTok{"Incumbent spend"}\NormalTok{),}
\NormalTok{  new\_m   }\SpecialCharTok{\%\textgreater{}\%} \FunctionTok{select}\NormalTok{(merchant\_id) }\SpecialCharTok{\%\textgreater{}\%} \FunctionTok{mutate}\NormalTok{(}\AttributeTok{window =} \StringTok{"New{-}merchant spend"}\NormalTok{)) }\SpecialCharTok{\%\textgreater{}\%}
  \FunctionTok{left\_join}\NormalTok{(tier, }\AttributeTok{by =} \StringTok{"merchant\_id"}\NormalTok{) }\SpecialCharTok{\%\textgreater{}\%}
  \FunctionTok{filter}\NormalTok{(}\SpecialCharTok{!}\FunctionTok{is.na}\NormalTok{(most\_recent\_sales\_range)) }\SpecialCharTok{\%\textgreater{}\%}
  \FunctionTok{count}\NormalTok{(window, most\_recent\_sales\_range) }\SpecialCharTok{\%\textgreater{}\%}
  \FunctionTok{group\_by}\NormalTok{(window) }\SpecialCharTok{\%\textgreater{}\%} \FunctionTok{mutate}\NormalTok{(}\AttributeTok{share =}\NormalTok{ n }\SpecialCharTok{/} \FunctionTok{sum}\NormalTok{(n)) }\SpecialCharTok{\%\textgreater{}\%} \FunctionTok{ungroup}\NormalTok{() }\SpecialCharTok{\%\textgreater{}\%}
  \FunctionTok{bind\_rows}\NormalTok{(size\_pyramid }\SpecialCharTok{\%\textgreater{}\%} \FunctionTok{transmute}\NormalTok{(}\AttributeTok{window =} \StringTok{"Merchant population"}\NormalTok{,}
\NormalTok{                                       most\_recent\_sales\_range, n, share))}

\NormalTok{tier\_flow }\SpecialCharTok{\%\textgreater{}\%} \FunctionTok{select}\NormalTok{(}\SpecialCharTok{{-}}\NormalTok{n) }\SpecialCharTok{\%\textgreater{}\%}
  \FunctionTok{pivot\_wider}\NormalTok{(}\AttributeTok{names\_from =}\NormalTok{ window, }\AttributeTok{values\_from =}\NormalTok{ share) }\SpecialCharTok{\%\textgreater{}\%}
  \FunctionTok{mutate}\NormalTok{(}\FunctionTok{across}\NormalTok{(}\FunctionTok{where}\NormalTok{(is.numeric), }\SpecialCharTok{\textasciitilde{}}\FunctionTok{percent}\NormalTok{(.x, }\FloatTok{0.1}\NormalTok{))) }\SpecialCharTok{\%\textgreater{}\%}
  \FunctionTok{kable}\NormalTok{(}\AttributeTok{col.names =} \FunctionTok{c}\NormalTok{(}\StringTok{"Tier"}\NormalTok{, }\StringTok{"Incumbent spend"}\NormalTok{, }\StringTok{"New{-}merchant spend"}\NormalTok{, }\StringTok{"Merchant population"}\NormalTok{))}
\end{Highlighting}
\end{Shaded}

{\def\LTcaptype{none} % do not increment counter
\begin{longtable}[]{@{}llll@{}}
\toprule\noalign{}
Tier & Incumbent spend & New-merchant spend & Merchant population \\
\midrule\noalign{}
\endhead
\bottomrule\noalign{}
\endlastfoot
A & 26.1\% & 16.9\% & 0.3\% \\
B & 19.1\% & 16.6\% & 1.5\% \\
C & 24.8\% & 27.9\% & 10.2\% \\
D & 19.1\% & 26.8\% & 35.1\% \\
E & 11.0\% & 11.9\% & 52.9\% \\
\end{longtable}
}

\begin{Shaded}
\begin{Highlighting}[]
\FunctionTok{ggplot}\NormalTok{(tier\_flow, }\FunctionTok{aes}\NormalTok{(most\_recent\_sales\_range, share,}
                      \AttributeTok{fill =} \FunctionTok{factor}\NormalTok{(window, }\AttributeTok{levels =} \FunctionTok{c}\NormalTok{(}\StringTok{"Merchant population"}\NormalTok{,}
                                                       \StringTok{"Incumbent spend"}\NormalTok{, }\StringTok{"New{-}merchant spend"}\NormalTok{)))) }\SpecialCharTok{+}
  \FunctionTok{geom\_col}\NormalTok{(}\AttributeTok{position =} \StringTok{"dodge"}\NormalTok{) }\SpecialCharTok{+}
  \FunctionTok{scale\_y\_continuous}\NormalTok{(}\AttributeTok{labels =}\NormalTok{ percent) }\SpecialCharTok{+}
  \FunctionTok{scale\_fill\_manual}\NormalTok{(}\AttributeTok{values =}\NormalTok{ PAL[}\FunctionTok{c}\NormalTok{(}\DecValTok{6}\NormalTok{, }\DecValTok{1}\NormalTok{, }\DecValTok{4}\NormalTok{)], }\AttributeTok{name =} \ConstantTok{NULL}\NormalTok{) }\SpecialCharTok{+}
  \FunctionTok{labs}\NormalTok{(}\AttributeTok{title =} \StringTok{"Incumbency sits at the top of the pyramid, discovery in the tail"}\NormalTok{,}
       \AttributeTok{subtitle =} \StringTok{"Tier A is 0.3\% of merchants and 26\% of incumbent volume {-} but only 17\% of new{-}merchant volume"}\NormalTok{,}
       \AttributeTok{x =} \StringTok{"Merchant size tier (A = largest)"}\NormalTok{, }\AttributeTok{y =} \StringTok{"Share"}\NormalTok{)}
\end{Highlighting}
\end{Shaded}

\begin{center}\includegraphics{Elo_Analysis_files/figure-latex/tiers-1} \end{center}

This is the most commercially loaded chart in the report.

\begin{itemize}
\tightlist
\item
  Tier A is \textbf{0.3\% of merchants} but takes \textbf{26.1\%} of
  incumbent transactions. These are the merchants a cardholder already
  uses.
\item
  In the novelty window that same tier takes only \textbf{16.9\%}, while
  tier D rises from 19.1\% to \textbf{26.8\%}.
\end{itemize}

\textbf{New commercial relationships form disproportionately in the long
tail.} A promotion placed with a tier A merchant is, on this evidence,
largely paying for traffic that was going to happen anyway. A promotion
placed with a healthy tier C or D merchant is far more likely to be
paying for a relationship that did not previously exist.

\begin{center}\rule{0.5\linewidth}{0.5pt}\end{center}

\section{Region: where supply does not match
demand}\label{region-where-supply-does-not-match-demand}

\begin{Shaded}
\begin{Highlighting}[]
\NormalTok{geo }\OtherTok{\textless{}{-}} \FunctionTok{full\_join}\NormalTok{(hist\_ok }\SpecialCharTok{\%\textgreater{}\%} \FunctionTok{count}\NormalTok{(state\_id, }\AttributeTok{name =} \StringTok{"inc\_txns"}\NormalTok{),}
\NormalTok{                 new\_w   }\SpecialCharTok{\%\textgreater{}\%} \FunctionTok{count}\NormalTok{(state\_id, }\AttributeTok{name =} \StringTok{"nov\_txns"}\NormalTok{), }\AttributeTok{by =} \StringTok{"state\_id"}\NormalTok{) }\SpecialCharTok{\%\textgreater{}\%}
  \FunctionTok{full\_join}\NormalTok{(merch }\SpecialCharTok{\%\textgreater{}\%} \FunctionTok{count}\NormalTok{(state\_id, }\AttributeTok{name =} \StringTok{"merchants"}\NormalTok{), }\AttributeTok{by =} \StringTok{"state\_id"}\NormalTok{) }\SpecialCharTok{\%\textgreater{}\%}
  \FunctionTok{replace\_na}\NormalTok{(}\FunctionTok{list}\NormalTok{(}\AttributeTok{inc\_txns =} \DecValTok{0}\NormalTok{, }\AttributeTok{nov\_txns =} \DecValTok{0}\NormalTok{, }\AttributeTok{merchants =} \DecValTok{0}\NormalTok{)) }\SpecialCharTok{\%\textgreater{}\%}
  \FunctionTok{mutate}\NormalTok{(}\StringTok{\textasciigrave{}}\AttributeTok{Demand share}\StringTok{\textasciigrave{}}   \OtherTok{=}\NormalTok{ inc\_txns }\SpecialCharTok{/} \FunctionTok{sum}\NormalTok{(inc\_txns),}
         \StringTok{\textasciigrave{}}\AttributeTok{Novelty share}\StringTok{\textasciigrave{}}  \OtherTok{=}\NormalTok{ nov\_txns }\SpecialCharTok{/} \FunctionTok{sum}\NormalTok{(nov\_txns),}
         \StringTok{\textasciigrave{}}\AttributeTok{Supply share}\StringTok{\textasciigrave{}}   \OtherTok{=}\NormalTok{ merchants }\SpecialCharTok{/} \FunctionTok{sum}\NormalTok{(merchants),}
         \StringTok{\textasciigrave{}}\AttributeTok{Coverage index}\StringTok{\textasciigrave{}} \OtherTok{=} \StringTok{\textasciigrave{}}\AttributeTok{Supply share}\StringTok{\textasciigrave{}} \SpecialCharTok{/} \StringTok{\textasciigrave{}}\AttributeTok{Demand share}\StringTok{\textasciigrave{}}\NormalTok{,}
         \StringTok{\textasciigrave{}}\AttributeTok{Novelty index}\StringTok{\textasciigrave{}}  \OtherTok{=} \StringTok{\textasciigrave{}}\AttributeTok{Novelty share}\StringTok{\textasciigrave{}} \SpecialCharTok{/} \StringTok{\textasciigrave{}}\AttributeTok{Demand share}\StringTok{\textasciigrave{}}\NormalTok{) }\SpecialCharTok{\%\textgreater{}\%}
  \FunctionTok{arrange}\NormalTok{(}\FunctionTok{desc}\NormalTok{(inc\_txns))}

\NormalTok{geo }\SpecialCharTok{\%\textgreater{}\%} \FunctionTok{head}\NormalTok{(}\DecValTok{10}\NormalTok{) }\SpecialCharTok{\%\textgreater{}\%}
  \FunctionTok{mutate}\NormalTok{(}\FunctionTok{across}\NormalTok{(}\FunctionTok{c}\NormalTok{(}\StringTok{\textasciigrave{}}\AttributeTok{Demand share}\StringTok{\textasciigrave{}}\NormalTok{, }\StringTok{\textasciigrave{}}\AttributeTok{Novelty share}\StringTok{\textasciigrave{}}\NormalTok{, }\StringTok{\textasciigrave{}}\AttributeTok{Supply share}\StringTok{\textasciigrave{}}\NormalTok{), }\SpecialCharTok{\textasciitilde{}}\FunctionTok{percent}\NormalTok{(.x, }\FloatTok{0.1}\NormalTok{)),}
         \FunctionTok{across}\NormalTok{(}\FunctionTok{c}\NormalTok{(}\StringTok{\textasciigrave{}}\AttributeTok{Coverage index}\StringTok{\textasciigrave{}}\NormalTok{, }\StringTok{\textasciigrave{}}\AttributeTok{Novelty index}\StringTok{\textasciigrave{}}\NormalTok{), }\SpecialCharTok{\textasciitilde{}}\FunctionTok{round}\NormalTok{(.x, }\DecValTok{2}\NormalTok{))) }\SpecialCharTok{\%\textgreater{}\%}
  \FunctionTok{select}\NormalTok{(state\_id, inc\_txns, }\StringTok{\textasciigrave{}}\AttributeTok{Demand share}\StringTok{\textasciigrave{}}\NormalTok{, }\StringTok{\textasciigrave{}}\AttributeTok{Supply share}\StringTok{\textasciigrave{}}\NormalTok{, }\StringTok{\textasciigrave{}}\AttributeTok{Coverage index}\StringTok{\textasciigrave{}}\NormalTok{,}
         \StringTok{\textasciigrave{}}\AttributeTok{Novelty share}\StringTok{\textasciigrave{}}\NormalTok{, }\StringTok{\textasciigrave{}}\AttributeTok{Novelty index}\StringTok{\textasciigrave{}}\NormalTok{) }\SpecialCharTok{\%\textgreater{}\%}
  \FunctionTok{kable}\NormalTok{(}\AttributeTok{col.names =} \FunctionTok{c}\NormalTok{(}\StringTok{"State"}\NormalTok{, }\StringTok{"Transactions"}\NormalTok{, }\StringTok{"Demand share"}\NormalTok{, }\StringTok{"Supply share"}\NormalTok{,}
                      \StringTok{"Coverage"}\NormalTok{, }\StringTok{"Novelty share"}\NormalTok{, }\StringTok{"Novelty index"}\NormalTok{),}
        \AttributeTok{format.args =} \FunctionTok{list}\NormalTok{(}\AttributeTok{big.mark =} \StringTok{","}\NormalTok{))}
\end{Highlighting}
\end{Shaded}

{\def\LTcaptype{none} % do not increment counter
\begin{longtable}[]{@{}
  >{\raggedleft\arraybackslash}p{(\linewidth - 12\tabcolsep) * \real{0.0732}}
  >{\raggedleft\arraybackslash}p{(\linewidth - 12\tabcolsep) * \real{0.1585}}
  >{\raggedright\arraybackslash}p{(\linewidth - 12\tabcolsep) * \real{0.1585}}
  >{\raggedright\arraybackslash}p{(\linewidth - 12\tabcolsep) * \real{0.1585}}
  >{\raggedleft\arraybackslash}p{(\linewidth - 12\tabcolsep) * \real{0.1098}}
  >{\raggedright\arraybackslash}p{(\linewidth - 12\tabcolsep) * \real{0.1707}}
  >{\raggedleft\arraybackslash}p{(\linewidth - 12\tabcolsep) * \real{0.1707}}@{}}
\toprule\noalign{}
\begin{minipage}[b]{\linewidth}\raggedleft
State
\end{minipage} & \begin{minipage}[b]{\linewidth}\raggedleft
Transactions
\end{minipage} & \begin{minipage}[b]{\linewidth}\raggedright
Demand share
\end{minipage} & \begin{minipage}[b]{\linewidth}\raggedright
Supply share
\end{minipage} & \begin{minipage}[b]{\linewidth}\raggedleft
Coverage
\end{minipage} & \begin{minipage}[b]{\linewidth}\raggedright
Novelty share
\end{minipage} & \begin{minipage}[b]{\linewidth}\raggedleft
Novelty index
\end{minipage} \\
\midrule\noalign{}
\endhead
\bottomrule\noalign{}
\endlastfoot
9 & 36,105 & 38.9\% & 27.7\% & 0.71 & 39.2\% & 1.01 \\
16 & 7,272 & 7.8\% & 9.9\% & 1.27 & 6.5\% & 0.84 \\
-1 & 6,203 & 6.7\% & 3.5\% & 0.53 & 4.7\% & 0.70 \\
15 & 5,493 & 5.9\% & 8.1\% & 1.38 & 8.0\% & 1.35 \\
5 & 5,351 & 5.8\% & 7.0\% & 1.22 & 7.0\% & 1.21 \\
20 & 4,536 & 4.9\% & 4.7\% & 0.96 & 4.2\% & 0.85 \\
11 & 3,913 & 4.2\% & 2.2\% & 0.52 & 2.4\% & 0.57 \\
21 & 3,682 & 4.0\% & 4.1\% & 1.03 & 2.6\% & 0.64 \\
19 & 3,387 & 3.6\% & 5.6\% & 1.54 & 4.2\% & 1.15 \\
22 & 2,894 & 3.1\% & 2.0\% & 0.63 & 4.1\% & 1.31 \\
\end{longtable}
}

\begin{Shaded}
\begin{Highlighting}[]
\NormalTok{geo }\SpecialCharTok{\%\textgreater{}\%} \FunctionTok{filter}\NormalTok{(inc\_txns }\SpecialCharTok{\textgreater{}=} \DecValTok{500}\NormalTok{) }\SpecialCharTok{\%\textgreater{}\%}
  \FunctionTok{ggplot}\NormalTok{(}\FunctionTok{aes}\NormalTok{(}\StringTok{\textasciigrave{}}\AttributeTok{Demand share}\StringTok{\textasciigrave{}}\NormalTok{, }\StringTok{\textasciigrave{}}\AttributeTok{Supply share}\StringTok{\textasciigrave{}}\NormalTok{)) }\SpecialCharTok{+}
  \FunctionTok{geom\_abline}\NormalTok{(}\AttributeTok{linetype =} \StringTok{"dashed"}\NormalTok{, }\AttributeTok{colour =} \StringTok{"grey60"}\NormalTok{) }\SpecialCharTok{+}
  \FunctionTok{geom\_point}\NormalTok{(}\FunctionTok{aes}\NormalTok{(}\AttributeTok{size =}\NormalTok{ inc\_txns), }\AttributeTok{colour =}\NormalTok{ PAL[}\DecValTok{1}\NormalTok{], }\AttributeTok{alpha =}\NormalTok{ .}\DecValTok{75}\NormalTok{) }\SpecialCharTok{+}
  \FunctionTok{geom\_text}\NormalTok{(}\FunctionTok{aes}\NormalTok{(}\AttributeTok{label =} \FunctionTok{paste0}\NormalTok{(}\StringTok{"S"}\NormalTok{, state\_id)), }\AttributeTok{vjust =} \SpecialCharTok{{-}}\FloatTok{1.1}\NormalTok{, }\AttributeTok{size =} \DecValTok{3}\NormalTok{, }\AttributeTok{colour =} \StringTok{"grey30"}\NormalTok{) }\SpecialCharTok{+}
  \FunctionTok{scale\_x\_continuous}\NormalTok{(}\AttributeTok{labels =}\NormalTok{ percent) }\SpecialCharTok{+} \FunctionTok{scale\_y\_continuous}\NormalTok{(}\AttributeTok{labels =}\NormalTok{ percent) }\SpecialCharTok{+}
  \FunctionTok{scale\_size\_continuous}\NormalTok{(}\AttributeTok{range =} \FunctionTok{c}\NormalTok{(}\DecValTok{2}\NormalTok{, }\DecValTok{10}\NormalTok{), }\AttributeTok{guide =} \StringTok{"none"}\NormalTok{) }\SpecialCharTok{+}
  \FunctionTok{labs}\NormalTok{(}\AttributeTok{title =} \StringTok{"Where the merchant base does not match the cardholders"}\NormalTok{,}
       \AttributeTok{subtitle =} \StringTok{"Below the dashed line = more cardholder demand than merchant supply"}\NormalTok{,}
       \AttributeTok{x =} \StringTok{"Share of card transactions"}\NormalTok{, }\AttributeTok{y =} \StringTok{"Share of merchants"}\NormalTok{,}
       \AttributeTok{caption =} \StringTok{"States with at least 500 transactions in the historical window. S{-}1 = unknown/unmapped state code."}\NormalTok{)}
\end{Highlighting}
\end{Shaded}

\begin{center}\includegraphics{Elo_Analysis_files/figure-latex/geo-1} \end{center}

State 9 alone is \textbf{39\% of card demand} but only \textbf{28\% of
the merchant base} --- a coverage index of 0.71, the largest
under-served market in the network. States 15, 19 and 16 carry more
merchants than their share of demand.

The more interesting column is the \textbf{novelty index}. State 22 has
only 3.1\% of demand but a novelty index of \textbf{1.31} --- its
cardholders try new merchants far more readily than the network average,
on a merchant base that is proportionally thin. That combination is a
merchant-acquisition target, not a discount target.

\begin{center}\rule{0.5\linewidth}{0.5pt}\end{center}

\section{The novelty window: what discovery actually looks
like}\label{the-novelty-window-what-discovery-actually-looks-like}

This is the file most teams ignore, and it is the only place in the
export where a \emph{new} commercial relationship is visible. A
promotion that works creates exactly this event.

\begin{Shaded}
\begin{Highlighting}[]
\FunctionTok{data.frame}\NormalTok{(}
  \AttributeTok{Metric =} \FunctionTok{c}\NormalTok{(}\StringTok{"Transactions"}\NormalTok{, }\StringTok{"Cards"}\NormalTok{, }\StringTok{"Distinct merchants"}\NormalTok{,}
             \StringTok{"New merchants tried per card (2 months)"}\NormalTok{,}
             \StringTok{"Average ticket at a new merchant (BRL)"}\NormalTok{,}
             \StringTok{"Average ticket at a known merchant (BRL)"}\NormalTok{,}
             \StringTok{"Median ticket at a new merchant (BRL)"}\NormalTok{,}
             \StringTok{"Median ticket at a known merchant (BRL)"}\NormalTok{),}
  \AttributeTok{Value =} \FunctionTok{c}\NormalTok{(}\FunctionTok{comma}\NormalTok{(}\FunctionTok{nrow}\NormalTok{(new\_w)), }\FunctionTok{comma}\NormalTok{(}\FunctionTok{n\_distinct}\NormalTok{(new\_w}\SpecialCharTok{$}\NormalTok{card\_id)),}
            \FunctionTok{comma}\NormalTok{(}\FunctionTok{n\_distinct}\NormalTok{(new\_m}\SpecialCharTok{$}\NormalTok{merchant\_id)),}
            \FunctionTok{round}\NormalTok{(}\FunctionTok{mean}\NormalTok{(}\FunctionTok{tapply}\NormalTok{(new\_m}\SpecialCharTok{$}\NormalTok{merchant\_id, new\_m}\SpecialCharTok{$}\NormalTok{card\_id,}
                              \ControlFlowTok{function}\NormalTok{(x) }\FunctionTok{length}\NormalTok{(}\FunctionTok{unique}\NormalTok{(x)))), }\DecValTok{1}\NormalTok{),}
            \FunctionTok{round}\NormalTok{(}\FunctionTok{mean}\NormalTok{(new\_w}\SpecialCharTok{$}\NormalTok{amount\_brl), }\DecValTok{2}\NormalTok{), }\FunctionTok{round}\NormalTok{(}\FunctionTok{mean}\NormalTok{(hist\_ok}\SpecialCharTok{$}\NormalTok{amount\_brl), }\DecValTok{2}\NormalTok{),}
            \FunctionTok{round}\NormalTok{(}\FunctionTok{median}\NormalTok{(new\_w}\SpecialCharTok{$}\NormalTok{amount\_brl), }\DecValTok{2}\NormalTok{), }\FunctionTok{round}\NormalTok{(}\FunctionTok{median}\NormalTok{(hist\_ok}\SpecialCharTok{$}\NormalTok{amount\_brl), }\DecValTok{2}\NormalTok{))) }\SpecialCharTok{\%\textgreater{}\%}
  \FunctionTok{kable}\NormalTok{()}
\end{Highlighting}
\end{Shaded}

{\def\LTcaptype{none} % do not increment counter
\begin{longtable}[]{@{}ll@{}}
\toprule\noalign{}
Metric & Value \\
\midrule\noalign{}
\endhead
\bottomrule\noalign{}
\endlastfoot
Transactions & 99,900 \\
Cards & 7,527 \\
Distinct merchants & 48,941 \\
New merchants tried per card (2 months) & 13.1 \\
Average ticket at a new merchant (BRL) & 118.7 \\
Average ticket at a known merchant (BRL) & 84.7 \\
Median ticket at a new merchant (BRL) & 48.05 \\
Median ticket at a known merchant (BRL) & 34.62 \\
\end{longtable}
}

\begin{Shaded}
\begin{Highlighting}[]
\NormalTok{novelty\_instal }\OtherTok{\textless{}{-}} \FunctionTok{bind\_rows}\NormalTok{(}
\NormalTok{  hist\_ok }\SpecialCharTok{\%\textgreater{}\%} \FunctionTok{filter}\NormalTok{(}\SpecialCharTok{!}\FunctionTok{is.na}\NormalTok{(instal\_band)) }\SpecialCharTok{\%\textgreater{}\%} \FunctionTok{count}\NormalTok{(instal\_band) }\SpecialCharTok{\%\textgreater{}\%}
    \FunctionTok{mutate}\NormalTok{(}\AttributeTok{Window =} \StringTok{"At a known merchant"}\NormalTok{, }\AttributeTok{share =}\NormalTok{ n }\SpecialCharTok{/} \FunctionTok{sum}\NormalTok{(n)),}
\NormalTok{  new\_w }\SpecialCharTok{\%\textgreater{}\%} \FunctionTok{filter}\NormalTok{(}\SpecialCharTok{!}\FunctionTok{is.na}\NormalTok{(instal\_band)) }\SpecialCharTok{\%\textgreater{}\%} \FunctionTok{count}\NormalTok{(instal\_band) }\SpecialCharTok{\%\textgreater{}\%}
    \FunctionTok{mutate}\NormalTok{(}\AttributeTok{Window =} \StringTok{"At a new merchant"}\NormalTok{, }\AttributeTok{share =}\NormalTok{ n }\SpecialCharTok{/} \FunctionTok{sum}\NormalTok{(n)))}

\FunctionTok{ggplot}\NormalTok{(novelty\_instal, }\FunctionTok{aes}\NormalTok{(instal\_band, share, }\AttributeTok{fill =}\NormalTok{ Window)) }\SpecialCharTok{+}
  \FunctionTok{geom\_col}\NormalTok{(}\AttributeTok{position =} \StringTok{"dodge"}\NormalTok{) }\SpecialCharTok{+}
  \FunctionTok{scale\_y\_continuous}\NormalTok{(}\AttributeTok{labels =}\NormalTok{ percent) }\SpecialCharTok{+}
  \FunctionTok{scale\_fill\_manual}\NormalTok{(}\AttributeTok{values =}\NormalTok{ PAL[}\FunctionTok{c}\NormalTok{(}\DecValTok{4}\NormalTok{, }\DecValTok{1}\NormalTok{)], }\AttributeTok{name =} \ConstantTok{NULL}\NormalTok{) }\SpecialCharTok{+}
  \FunctionTok{labs}\NormalTok{(}\AttributeTok{title =} \StringTok{"Trying somewhere new is a considered purchase"}\NormalTok{,}
       \AttributeTok{subtitle =} \StringTok{"A first visit carries a bigger ticket and is more likely to be split into instalments"}\NormalTok{,}
       \AttributeTok{x =} \StringTok{"Instalments"}\NormalTok{, }\AttributeTok{y =} \StringTok{"Share of transactions"}\NormalTok{)}
\end{Highlighting}
\end{Shaded}

\begin{center}\includegraphics{Elo_Analysis_files/figure-latex/novelty-instal-1} \end{center}

Three findings, each one actionable:

\begin{enumerate}
\def\labelenumi{\arabic{enumi}.}
\tightlist
\item
  \textbf{A first visit costs 40\% more.} R\$118.70 at a new merchant
  against R\$84.70 at a known one. Discovery is not impulse --- it is a
  considered purchase, and it is worth more per transaction.
\item
  \textbf{Instalment use nearly doubles.} 2--3 instalment purchases run
  at 4.2\% at new merchants against 2.5\% at known ones. The one place
  where Elo's cardholders do reach for instalments is a first purchase
  somewhere new. That is where an instalment-linked promotion would
  land, and nowhere else.
\item
  \textbf{Cardholders try roughly 13 new merchants in two months.}
  Appetite for discovery is not the constraint. Direction is.
\end{enumerate}

\begin{center}\rule{0.5\linewidth}{0.5pt}\end{center}

\section{The Q4 targeting map}\label{the-q4-targeting-map}

We now put the demand side and supply side on one picture. Two indices
per merchant subsector:

\begin{itemize}
\tightlist
\item
  \textbf{Novelty Capture Index (NCI)} = share of new-merchant volume ÷
  share of incumbent volume. \textbf{Above 1} = cardholders come here to
  try something new. \textbf{Below 1} = they come here out of habit.
\item
  \textbf{Coverage Index (COV)} = share of merchants ÷ share of
  incumbent volume. \textbf{Above 1} = more merchant supply than
  cardholder demand. \textbf{Below 1} = demand outruns supply.
\end{itemize}

\begin{Shaded}
\begin{Highlighting}[]
\NormalTok{sub\_map }\OtherTok{\textless{}{-}} \FunctionTok{full\_join}\NormalTok{(hist\_ok }\SpecialCharTok{\%\textgreater{}\%} \FunctionTok{count}\NormalTok{(subsector\_id, }\AttributeTok{name =} \StringTok{"inc\_txns"}\NormalTok{),}
\NormalTok{                     new\_w   }\SpecialCharTok{\%\textgreater{}\%} \FunctionTok{count}\NormalTok{(subsector\_id, }\AttributeTok{name =} \StringTok{"nov\_txns"}\NormalTok{), }\AttributeTok{by =} \StringTok{"subsector\_id"}\NormalTok{) }\SpecialCharTok{\%\textgreater{}\%}
  \FunctionTok{full\_join}\NormalTok{(merch }\SpecialCharTok{\%\textgreater{}\%} \FunctionTok{count}\NormalTok{(subsector\_id, }\AttributeTok{name =} \StringTok{"merchants"}\NormalTok{), }\AttributeTok{by =} \StringTok{"subsector\_id"}\NormalTok{) }\SpecialCharTok{\%\textgreater{}\%}
  \FunctionTok{full\_join}\NormalTok{(hist\_ok }\SpecialCharTok{\%\textgreater{}\%} \FunctionTok{group\_by}\NormalTok{(subsector\_id) }\SpecialCharTok{\%\textgreater{}\%}
              \FunctionTok{summarise}\NormalTok{(}\AttributeTok{avg\_ticket =} \FunctionTok{mean}\NormalTok{(amount\_brl), }\AttributeTok{.groups =} \StringTok{"drop"}\NormalTok{), }\AttributeTok{by =} \StringTok{"subsector\_id"}\NormalTok{) }\SpecialCharTok{\%\textgreater{}\%}
  \FunctionTok{replace\_na}\NormalTok{(}\FunctionTok{list}\NormalTok{(}\AttributeTok{inc\_txns =} \DecValTok{0}\NormalTok{, }\AttributeTok{nov\_txns =} \DecValTok{0}\NormalTok{, }\AttributeTok{merchants =} \DecValTok{0}\NormalTok{)) }\SpecialCharTok{\%\textgreater{}\%}
  \FunctionTok{mutate}\NormalTok{(}\AttributeTok{demand\_share =}\NormalTok{ inc\_txns }\SpecialCharTok{/} \FunctionTok{sum}\NormalTok{(inc\_txns),}
         \AttributeTok{novelty\_share =}\NormalTok{ nov\_txns }\SpecialCharTok{/} \FunctionTok{sum}\NormalTok{(nov\_txns),}
         \AttributeTok{supply\_share =}\NormalTok{ merchants }\SpecialCharTok{/} \FunctionTok{sum}\NormalTok{(merchants),}
         \AttributeTok{NCI =}\NormalTok{ novelty\_share }\SpecialCharTok{/}\NormalTok{ demand\_share,}
         \AttributeTok{COV =}\NormalTok{ supply\_share }\SpecialCharTok{/}\NormalTok{ demand\_share) }\SpecialCharTok{\%\textgreater{}\%}
  \FunctionTok{filter}\NormalTok{(inc\_txns }\SpecialCharTok{\textgreater{}=} \DecValTok{300}\NormalTok{) }\SpecialCharTok{\%\textgreater{}\%}
  \FunctionTok{mutate}\NormalTok{(}\AttributeTok{Role =} \FunctionTok{case\_when}\NormalTok{(}
\NormalTok{    NCI }\SpecialCharTok{\textgreater{}=} \FloatTok{1.2}              \SpecialCharTok{\textasciitilde{}} \StringTok{"Acquisition engine"}\NormalTok{,}
\NormalTok{    COV }\SpecialCharTok{\textgreater{}=} \FloatTok{1.2}              \SpecialCharTok{\textasciitilde{}} \StringTok{"Over{-}served {-} promotion waste"}\NormalTok{,}
    \ConstantTok{TRUE}                    \SpecialCharTok{\textasciitilde{}} \StringTok{"Habit anchor {-} under{-}supplied"}\NormalTok{)) }\SpecialCharTok{\%\textgreater{}\%}
  \FunctionTok{arrange}\NormalTok{(}\FunctionTok{desc}\NormalTok{(demand\_share))}

\FunctionTok{ggplot}\NormalTok{(sub\_map, }\FunctionTok{aes}\NormalTok{(COV, NCI)) }\SpecialCharTok{+}
  \FunctionTok{geom\_hline}\NormalTok{(}\AttributeTok{yintercept =} \FloatTok{1.2}\NormalTok{, }\AttributeTok{linetype =} \StringTok{"dashed"}\NormalTok{, }\AttributeTok{colour =} \StringTok{"grey55"}\NormalTok{) }\SpecialCharTok{+}
  \FunctionTok{geom\_vline}\NormalTok{(}\AttributeTok{xintercept =} \FloatTok{1.2}\NormalTok{, }\AttributeTok{linetype =} \StringTok{"dashed"}\NormalTok{, }\AttributeTok{colour =} \StringTok{"grey55"}\NormalTok{) }\SpecialCharTok{+}
  \FunctionTok{geom\_point}\NormalTok{(}\FunctionTok{aes}\NormalTok{(}\AttributeTok{size =}\NormalTok{ demand\_share, }\AttributeTok{colour =}\NormalTok{ Role), }\AttributeTok{alpha =}\NormalTok{ .}\DecValTok{85}\NormalTok{) }\SpecialCharTok{+}
  \FunctionTok{geom\_text}\NormalTok{(}\FunctionTok{aes}\NormalTok{(}\AttributeTok{label =}\NormalTok{ subsector\_id), }\AttributeTok{size =} \FloatTok{2.8}\NormalTok{, }\AttributeTok{vjust =} \SpecialCharTok{{-}}\FloatTok{1.2}\NormalTok{, }\AttributeTok{colour =} \StringTok{"grey25"}\NormalTok{) }\SpecialCharTok{+}
  \FunctionTok{scale\_size\_continuous}\NormalTok{(}\AttributeTok{range =} \FunctionTok{c}\NormalTok{(}\FloatTok{2.5}\NormalTok{, }\DecValTok{12}\NormalTok{), }\AttributeTok{labels =}\NormalTok{ percent, }\AttributeTok{name =} \StringTok{"Share of volume"}\NormalTok{) }\SpecialCharTok{+}
  \FunctionTok{scale\_colour\_manual}\NormalTok{(}\AttributeTok{values =} \FunctionTok{c}\NormalTok{(}\StringTok{"Acquisition engine"} \OtherTok{=}\NormalTok{ PAL[}\DecValTok{1}\NormalTok{],}
                                 \StringTok{"Habit anchor {-} under{-}supplied"} \OtherTok{=}\NormalTok{ PAL[}\DecValTok{4}\NormalTok{],}
                                 \StringTok{"Over{-}served {-} promotion waste"} \OtherTok{=}\NormalTok{ PAL[}\DecValTok{2}\NormalTok{]), }\AttributeTok{name =} \ConstantTok{NULL}\NormalTok{) }\SpecialCharTok{+}
  \FunctionTok{guides}\NormalTok{(}\AttributeTok{colour =} \FunctionTok{guide\_legend}\NormalTok{(}\AttributeTok{nrow =} \DecValTok{1}\NormalTok{, }\AttributeTok{order =} \DecValTok{1}\NormalTok{, }\AttributeTok{override.aes =} \FunctionTok{list}\NormalTok{(}\AttributeTok{size =} \DecValTok{4}\NormalTok{))) }\SpecialCharTok{+}
  \FunctionTok{labs}\NormalTok{(}\AttributeTok{title =} \StringTok{"Where to point the Q4 promotions programme"}\NormalTok{,}
       \AttributeTok{subtitle =} \StringTok{"Novelty Capture Index against Coverage Index, by merchant subsector"}\NormalTok{,}
       \AttributeTok{x =} \StringTok{"Coverage Index  (merchant supply / cardholder demand)"}\NormalTok{,}
       \AttributeTok{y =} \StringTok{"Novelty Capture Index  (new{-}merchant volume / incumbent volume)"}\NormalTok{,}
       \AttributeTok{caption =} \StringTok{"Bubble size = share of incumbent transactions. Subsectors with 300+ incumbent transactions."}\NormalTok{)}
\end{Highlighting}
\end{Shaded}

\begin{center}\includegraphics{Elo_Analysis_files/figure-latex/map-1} \end{center}

\begin{Shaded}
\begin{Highlighting}[]
\NormalTok{sub\_map }\SpecialCharTok{\%\textgreater{}\%}
  \FunctionTok{transmute}\NormalTok{(}\AttributeTok{Subsector =}\NormalTok{ subsector\_id, }\AttributeTok{Transactions =}\NormalTok{ inc\_txns,}
            \StringTok{\textasciigrave{}}\AttributeTok{Demand \%}\StringTok{\textasciigrave{}} \OtherTok{=} \FunctionTok{percent}\NormalTok{(demand\_share, }\FloatTok{0.1}\NormalTok{),}
            \StringTok{\textasciigrave{}}\AttributeTok{Novelty \%}\StringTok{\textasciigrave{}} \OtherTok{=} \FunctionTok{percent}\NormalTok{(novelty\_share, }\FloatTok{0.1}\NormalTok{),}
            \StringTok{\textasciigrave{}}\AttributeTok{Supply \%}\StringTok{\textasciigrave{}} \OtherTok{=} \FunctionTok{percent}\NormalTok{(supply\_share, }\FloatTok{0.1}\NormalTok{),}
            \AttributeTok{NCI =} \FunctionTok{round}\NormalTok{(NCI, }\DecValTok{2}\NormalTok{), }\AttributeTok{COV =} \FunctionTok{round}\NormalTok{(COV, }\DecValTok{2}\NormalTok{),}
            \StringTok{\textasciigrave{}}\AttributeTok{Avg ticket}\StringTok{\textasciigrave{}} \OtherTok{=} \FunctionTok{round}\NormalTok{(avg\_ticket, }\DecValTok{0}\NormalTok{), Role) }\SpecialCharTok{\%\textgreater{}\%}
  \FunctionTok{kable}\NormalTok{(}\AttributeTok{format.args =} \FunctionTok{list}\NormalTok{(}\AttributeTok{big.mark =} \StringTok{","}\NormalTok{))}
\end{Highlighting}
\end{Shaded}

{\def\LTcaptype{none} % do not increment counter
\begin{longtable}[]{@{}
  >{\raggedleft\arraybackslash}p{(\linewidth - 16\tabcolsep) * \real{0.0980}}
  >{\raggedleft\arraybackslash}p{(\linewidth - 16\tabcolsep) * \real{0.1275}}
  >{\raggedright\arraybackslash}p{(\linewidth - 16\tabcolsep) * \real{0.0882}}
  >{\raggedright\arraybackslash}p{(\linewidth - 16\tabcolsep) * \real{0.0980}}
  >{\raggedright\arraybackslash}p{(\linewidth - 16\tabcolsep) * \real{0.0882}}
  >{\raggedleft\arraybackslash}p{(\linewidth - 16\tabcolsep) * \real{0.0490}}
  >{\raggedleft\arraybackslash}p{(\linewidth - 16\tabcolsep) * \real{0.0490}}
  >{\raggedleft\arraybackslash}p{(\linewidth - 16\tabcolsep) * \real{0.1078}}
  >{\raggedright\arraybackslash}p{(\linewidth - 16\tabcolsep) * \real{0.2941}}@{}}
\toprule\noalign{}
\begin{minipage}[b]{\linewidth}\raggedleft
Subsector
\end{minipage} & \begin{minipage}[b]{\linewidth}\raggedleft
Transactions
\end{minipage} & \begin{minipage}[b]{\linewidth}\raggedright
Demand \%
\end{minipage} & \begin{minipage}[b]{\linewidth}\raggedright
Novelty \%
\end{minipage} & \begin{minipage}[b]{\linewidth}\raggedright
Supply \%
\end{minipage} & \begin{minipage}[b]{\linewidth}\raggedleft
NCI
\end{minipage} & \begin{minipage}[b]{\linewidth}\raggedleft
COV
\end{minipage} & \begin{minipage}[b]{\linewidth}\raggedleft
Avg ticket
\end{minipage} & \begin{minipage}[b]{\linewidth}\raggedright
Role
\end{minipage} \\
\midrule\noalign{}
\endhead
\bottomrule\noalign{}
\endlastfoot
37 & 17,486 & 18.8\% & 21.0\% & 10.6\% & 1.12 & 0.56 & 49 & Habit anchor
- under-supplied \\
33 & 17,246 & 18.6\% & 7.8\% & 9.2\% & 0.42 & 0.50 & 69 & Habit anchor -
under-supplied \\
19 & 11,161 & 12.0\% & 10.6\% & 6.1\% & 0.88 & 0.51 & 71 & Habit anchor
- under-supplied \\
34 & 10,507 & 11.3\% & 9.8\% & 7.9\% & 0.86 & 0.70 & 39 & Habit anchor -
under-supplied \\
29 & 8,269 & 8.9\% & 6.3\% & 6.4\% & 0.71 & 0.72 & 87 & Habit anchor -
under-supplied \\
16 & 4,962 & 5.3\% & 5.3\% & 4.9\% & 0.99 & 0.92 & 62 & Habit anchor -
under-supplied \\
27 & 4,111 & 4.4\% & 9.6\% & 15.4\% & 2.16 & 3.49 & 175 & Acquisition
engine \\
21 & 2,494 & 2.7\% & 3.4\% & 7.2\% & 1.26 & 2.70 & 148 & Acquisition
engine \\
7 & 2,445 & 2.6\% & 2.5\% & 2.3\% & 0.95 & 0.87 & 85 & Habit anchor -
under-supplied \\
1 & 1,845 & 2.0\% & 2.9\% & 1.8\% & 1.48 & 0.92 & 155 & Acquisition
engine \\
15 & 1,789 & 1.9\% & 1.3\% & 1.0\% & 0.68 & 0.53 & 183 & Habit anchor -
under-supplied \\
39 & 1,111 & 1.2\% & 1.4\% & 0.2\% & 1.20 & 0.21 & 20 & Habit anchor -
under-supplied \\
36 & 959 & 1.0\% & 1.9\% & 1.5\% & 1.81 & 1.41 & 98 & Acquisition
engine \\
18 & 865 & 0.9\% & 2.0\% & 4.3\% & 2.13 & 4.58 & 249 & Acquisition
engine \\
25 & 855 & 0.9\% & 1.4\% & 1.4\% & 1.48 & 1.47 & 104 & Acquisition
engine \\
30 & 852 & 0.9\% & 1.5\% & 1.3\% & 1.62 & 1.43 & 115 & Acquisition
engine \\
12 & 820 & 0.9\% & 1.5\% & 1.3\% & 1.66 & 1.53 & 180 & Acquisition
engine \\
20 & 707 & 0.8\% & 1.9\% & 2.8\% & 2.44 & 3.74 & 240 & Acquisition
engine \\
41 & 642 & 0.7\% & 0.7\% & 0.9\% & 0.96 & 1.33 & 110 & Over-served -
promotion waste \\
8 & 626 & 0.7\% & 0.5\% & 0.4\% & 0.81 & 0.61 & 42 & Habit anchor -
under-supplied \\
31 & 445 & 0.5\% & 0.8\% & 0.7\% & 1.66 & 1.50 & 93 & Acquisition
engine \\
17 & 383 & 0.4\% & 1.0\% & 0.1\% & 2.47 & 0.36 & 146 & Acquisition
engine \\
32 & 325 & 0.3\% & 0.8\% & 2.0\% & 2.18 & 5.72 & 315 & Acquisition
engine \\
\end{longtable}
}

\textbf{Reading the map.}

\emph{Top half --- acquisition engines (subsectors 27, 18, 20, 32, 36,
12, 30, 25, 31, 21, 17, 1).} Subsector 27 is the clearest case: 4.4\% of
incumbent volume but \textbf{9.6\% of novelty volume (NCI 2.16)},
sitting on 15.4\% of the merchant base (COV 3.49). Cardholders are
already willing to try new merchants here, and there is ample merchant
supply to promote with. Average ticket R\$175. This is where a discovery
offer converts, and where the supply side can absorb the traffic.

\emph{Bottom left --- habit anchors, under-supplied (33, 29, 19, 34, 37,
16, 15, 7, 8, 39).} Subsector 33 is the network's second largest at
18.6\% of volume, but an \textbf{NCI of 0.42} --- cardholders go where
they always go. Discounting here buys transactions that were happening
anyway. Its coverage index of 0.50 says demand outruns merchant supply,
so the right intervention is \textbf{merchant recruitment, not
cardholder discounting.}

\emph{Bottom right --- over-served (subsector 41).} More merchant supply
than either demand or discovery justifies. A small quadrant on this
data, but it is the one place where promotion spend is waste by
construction.

\emph{Subsector 37 deserves a line of its own.} At 18.8\% of
transactions it is the single largest, with an NCI of 1.12 --- just
below the discovery threshold --- and a R\$49 average ticket. It is the
workhorse of the network: the place to defend top-of-wallet position,
not the place to chase acquisition.

\begin{center}\rule{0.5\linewidth}{0.5pt}\end{center}

\section{Recommendations for Q4}\label{recommendations-for-q4}

\textbf{1. Stop discounting habit. Redirect that budget to discovery.}
Subsectors 33, 29, 19 and 34 together are \textbf{50.8\% of
transactions} and have novelty capture indices below 0.9. Money spent on
offers here is subsidising purchases that would have happened. Moving
even a third of that budget into the acquisition-engine quadrant (27,
18, 20, 32, 36) targets the same cardholders at the moment they are
actually willing to change behaviour.

\textbf{2. Promote down the pyramid, not up it.} Tier A merchants are
0.3\% of the base and 26\% of incumbent volume, but only 17\% of
new-merchant volume. Set a floor: a defined share of Q4 promotional
slots reserved for \textbf{accelerating tier C and D merchants}, which
is where new relationships actually form, and where a merchant will
value the traffic enough to fund part of the offer.

\textbf{3. Fix the \texttt{category\_1\ =\ Y} decline problem before
spending.} That flag carries \textbf{4.3× the odds of a decline} and
covers 6.1\% of volume. Promoting into a population we then turn away
costs money twice. This is a one-week question for the risk team and it
should be answered before the Q4 calendar is locked.

\textbf{4. Use instalments in exactly one place: the first visit.}
Instalment use nearly doubles at new merchants. An instalment-linked
offer at a \emph{known} merchant is fighting the network's nature; at a
first visit it works with it. This is the one instalment play the data
supports.

\textbf{5. Recruit merchants where demand outruns supply.} State 9 (39\%
of demand, 28\% of supply) and subsector 33 (18.6\% of demand, coverage
index 0.50) are merchant-acquisition problems wearing a promotions
costume. State 22 --- high novelty appetite, thin supply --- is the
cleanest greenfield in the network.

\textbf{6. Segment the send.} Core Everyday (72\% of spend) needs
\textbf{retention and top-of-wallet defence}, not discounts. Big-Ticket
Narrow (43 days since last transaction) needs \textbf{reactivation}.
Explorer Light is the natural \textbf{discovery audience}. Occasional
--- a third of cards and 8.7\% of spend --- should be the last call on
the budget, not the first.

\begin{center}\rule{0.5\linewidth}{0.5pt}\end{center}

\section{What this analysis cannot tell
you}\label{what-this-analysis-cannot-tell-you}

Stating the limits is part of the answer.

\begin{enumerate}
\def\labelenumi{\arabic{enumi}.}
\tightlist
\item
  \textbf{No card-level loyalty file.} \texttt{train.csv} --- with
  \texttt{first\_active\_month} and the loyalty \texttt{target} --- was
  not in the export. We rebuilt tenure and value from observed
  behaviour, so we can describe what loyal behaviour \emph{looks like}
  but cannot explain the score the business already computes.
\item
  \textbf{The two windows are different card samples.} 16 cards in
  common, so every novelty finding is a portfolio-level comparison, not
  a within-card before/after. It supports ``where discovery happens'';
  it does not support ``this promotion caused this card to switch''.
\item
  \textbf{463 cards is a small panel.} Portfolio-level conclusions on
  instalment mix, ticket size and concentration are stable;
  state-by-subsector cells with a few hundred transactions are
  indicative only.
\item
  \textbf{Everything is anonymised.} We can say subsector 27 behaves
  like a discovery category and subsector 33 like a habit category.
  Someone inside Elo has to put the names on them before this becomes a
  media plan.
\item
  \textbf{The currency transform is empirically recovered, not
  supplied.} It is validated by the shape of the resulting distribution
  --- no negative amounts, an economically coherent median --- but it
  should be confirmed against one known settlement record before any
  figure is quoted externally.
\item
  \textbf{Correlation, not causation.} The targeting map ranks where
  discovery already happens. Proving a promotion \emph{causes} it needs
  a holdout test, which is the natural next piece of work.
\end{enumerate}

\begin{center}\rule{0.5\linewidth}{0.5pt}\end{center}

\section{Major learnings}\label{major-learnings}

\begin{itemize}
\tightlist
\item
  \textbf{Audit before you analyse.} 90\% of the largest file was empty
  padding and a hundred records were about to hijack every average.
  Neither was flagged. The first hour of data work protected every
  number that followed.
\item
  \textbf{Make the data speak the Board's language.} Recovering the
  currency transform changed nothing statistically and changed
  everything about whether the analysis could be presented. ``The median
  ticket is −0.69'' is not a finding; ``the median ticket is R\$35'' is
  a strategy.
\item
  \textbf{The identity question was answered by a simple frequency
  table.} 96\% single-payment settled what kind of network Elo is. No
  model was required --- only asking the right question of the right
  column.
\item
  \textbf{The file everyone ignores was the one that mattered.} The
  novelty window is the only view of a \emph{new} relationship forming,
  which is precisely what a promotion is buying. Analysing it against
  the incumbent window turned a description into a recommendation.
\item
  \textbf{Ratios travel further than totals.} Raw volume said ``promote
  subsector 33, it is huge''. Volume \emph{relative to} novelty and
  supply said the opposite. Two constructed indices changed the answer.
\item
  \textbf{State the limits in the room.} Naming the 16-card overlap up
  front made the rest of the analysis credible rather than fragile.
\end{itemize}

\begin{verbatim}
## Analysis run: 28 Sep 2026 10:32 | R version 4.6.1 (2026-06-24)
\end{verbatim}

\end{document}
